\documentclass[12pt, titlepage]{article}
\pdfinfoomitdate=1
\pdftrailerid{}

\usepackage{booktabs}
\usepackage{tabularx}
\usepackage{longtable}
\usepackage{hyperref}
\usepackage{xcolor}
\hypersetup{
    colorlinks,
    citecolor=blue,
    filecolor=black,
    linkcolor=red,
    urlcolor=blue
}
\usepackage[round]{natbib}
\usepackage{pifont}
\usepackage{booktabs,tabularx}
\usepackage{float}
\usepackage{array}      % better column types
\usepackage{float}      % [H] placement
\usepackage{amssymb} 

\input{../Comments}
\input{../Common}
\newcommand{\mref}[1]{\texttt{#1}}

\newcommand{\tidanchor}[1]{\hypertarget{tid:#1}{}}
\newcommand{\tid}[1]{\hyperlink{tid:#1}{\textcolor{blue}{\texttt{#1}}}}
\newcommand{\tiddef}[1]{\textcolor{blue}{\texttt{#1}}}
\newcommand{\tidallanchors}{%
  \tidanchor{test-acquire-frame}%
  \tidanchor{test-api-client-automatic-discovery}%
  \tidanchor{test-api-client-methods}%
  \tidanchor{test-api-client-urls}%
  \tidanchor{test-api-endpoints-exist}%
  \tidanchor{test-api-error-responses}%
  \tidanchor{test-api-recordings-crud}%
  \tidanchor{test-api-validation}%
  \tidanchor{test-app-routing}%
  \tidanchor{test-bilingual-ui}%
  \tidanchor{test-bounding-box-center}%
  \tidanchor{test-bounding-box-collection}%
  \tidanchor{test-bounding-box-rotate-90}%
  \tidanchor{test-color-vision-deficiency}%
  \tidanchor{test-configuration-menu}%
  \tidanchor{test-confirmation-prompts}%
  \tidanchor{test-consistent-units}%
  \tidanchor{test-control-page-bbox-overlay}%
  \tidanchor{test-control-page-drag}%
  \tidanchor{test-control-page-render}%
  \tidanchor{test-control-process-main-entry}%
  \tidanchor{test-controls-card}%
  \tidanchor{test-controls-card-render}%
  \tidanchor{test-cv-format-time}%
  \tidanchor{test-cv-management-exception-safety}%
  \tidanchor{test-cv-management-osd-send}%
  \tidanchor{test-cv-management-pause-resume}%
  \tidanchor{test-cv-management-recording-commands}%
  \tidanchor{test-cv-management-recv-loop}%
  \tidanchor{test-cv-update-osd}%
  \tidanchor{test-database-allocate}%
  \tidanchor{test-database-error-handling}%
  \tidanchor{test-database-list-recordings}%
  \tidanchor{test-database-path-traversal-protection}%
  \tidanchor{test-database-rename-delete}%
  \tidanchor{test-database-space-usage}%
  \tidanchor{test-dataclass-construction}%
  \tidanchor{test-disconnect-camera}%
  \tidanchor{test-disconnect-hdmi}%
  \tidanchor{test-disk-full}%
  \tidanchor{test-exit-gimbal-range}%
  \tidanchor{test-formatters-bytes}%
  \tidanchor{test-formatters-degrees}%
  \tidanchor{test-formatters-duration}%
  \tidanchor{test-formatters-edge-cases}%
  \tidanchor{test-formatters-internationalization}%
  \tidanchor{test-formatters-temperature}%
  \tidanchor{test-gimbal-angle-commands}%
  \tidanchor{test-gimbal-crc8}%
  \tidanchor{test-gimbal-info-query}%
  \tidanchor{test-gimbal-interface}%
  \tidanchor{test-gimbal-led-control}%
  \tidanchor{test-gimbal-packet-construction}%
  \tidanchor{test-hdmi-output}%
  \tidanchor{test-i18n-activation}%
  \tidanchor{test-immediate-ui-feedback}%
  \tidanchor{test-ip4-addresses}%
  \tidanchor{test-ipc-buffer-memory-safety}%
  \tidanchor{test-ipc-buffer-overrun-handling}%
  \tidanchor{test-ipc-buffer-receiver}%
  \tidanchor{test-ipc-buffer-sender}%
  \tidanchor{test-language-selector}%
  \tidanchor{test-livestream-construction}%
  \tidanchor{test-livestream-process-loop}%
  \tidanchor{test-log-operations}%
  \tidanchor{test-main-dispatch-control}%
  \tidanchor{test-main-dispatch-cv}%
  \tidanchor{test-main-dispatch-livestream}%
  \tidanchor{test-main-dispatch-transcode}%
  \tidanchor{test-main-dispatch-unknown}%
  \tidanchor{test-manage-recordings}%
  \tidanchor{test-manual-arm-mode}%
  \tidanchor{test-manual-control}%
  \tidanchor{test-manual-idle-mode}%
  \tidanchor{test-navbar}%
  \tidanchor{test-network-disconnect}%
  \tidanchor{test-no-collection-pii}%
  \tidanchor{test-preview-real-time}%
  \tidanchor{test-recording}%
  \tidanchor{test-recordings-delete}%
  \tidanchor{test-recordings-page-empty}%
  \tidanchor{test-recordings-page-list}%
  \tidanchor{test-recordings-page-loading}%
  \tidanchor{test-recordings-rename}%
  \tidanchor{test-report-errors}%
  \tidanchor{test-return-to-idle}%
  \tidanchor{test-rocam-provider-error-handling}%
  \tidanchor{test-rocam-provider-initialization}%
  \tidanchor{test-rocam-provider-polling}%
  \tidanchor{test-scheduler-validation}%
  \tidanchor{test-sse-broadcast-performance}%
  \tidanchor{test-sse-thread-safety}%
  \tidanchor{test-state-error-handling}%
  \tidanchor{test-state-management-arm-disarm}%
  \tidanchor{test-state-management-download-pause}%
  \tidanchor{test-state-management-manual-move}%
  \tidanchor{test-state-management-manual-move-clamped}%
  \tidanchor{test-state-management-recording}%
  \tidanchor{test-state-management-recording-only-noops}%
  \tidanchor{test-state-management-recording-only-status}%
  \tidanchor{test-state-status-performance}%
  \tidanchor{test-stop-immediately}%
  \tidanchor{test-system-status-caching}%
  \tidanchor{test-system-status-callback-updates}%
  \tidanchor{test-system-status-card}%
  \tidanchor{test-tracking}%
  \tidanchor{test-tracking-angle-clamping}%
  \tidanchor{test-tracking-queue-behavior}%
  \tidanchor{test-tracking-stop}%
  \tidanchor{test-tracking-worker-control-law}%
  \tidanchor{test-transcode-bus-call}%
  \tidanchor{test-transcode-pipe-cleanup}%
  \tidanchor{test-transcode-read-log}%
  \tidanchor{test-transcode-routes-404}%
  \tidanchor{test-transcode-shader-probe}%
  \tidanchor{test-transcode-streaming}%
  \tidanchor{test-transition-to-tracking}%
  \tidanchor{test-ui-error-handling}%
  \tidanchor{test-ui-loading-states}%
  \tidanchor{test-use-rocam-hook}%
  \tidanchor{test-user-survey}%
  \tidanchor{test-validate-commands}%
  \tidanchor{test-visual-indicators}%
  \tidanchor{test-watchdog-clear}%
  \tidanchor{test-watchdog-refresh}%
  \tidanchor{test-watchdog-timeout-accuracy}%
}


\begin{document}
% Initialize hidden anchors for all Test IDs
\tidallanchors

% Links to project documents (relative to docs/VnVPlan/); cite keys in References.bib
\newcommand{\SRS}{\href{../SRS/SRS.pdf}{SRS}}
\newcommand{\MG}{\href{../Design/SoftArchitecture/MG.pdf}{MG}}
\newcommand{\MIS}{\href{../Design/SoftDetailedDes/MIS.pdf}{MIS}}

\title{System Verification and Validation Plan for \progname{}}
\author{\authname}
\date{March 9, 2026}

\maketitle

\pagenumbering{roman}

\section*{Revision History}

\begin{tabularx}{\textwidth}{p{3cm}p{2cm}X}
  \toprule {\bf Date} & {\bf Version} & {\bf Notes}                                                                                                                                                                                                                       \\
  \midrule
  2025-10-21          & 1.0           & Initial Draft                                                                                                                                                                                                                     \\
  2025-10-23          & 2.0           & Revised on section 3                                                                                                                                                                                                              \\
  2025-10-26          & 3.0           & Revised on section 3 based on TA and peer feedback                                                                                                                                                                                \\
  2025-10-27          & 4.0           & Added system tests                                                                                                                                                                                                                \\
  2026-03-04          & 5.0           & Scope limited to unit testing and stakeholder (McMaster Rocketry Team) review only; removed supervisor, classmate, task-based inspection, field testing                                                                           \\
  2026-03-05          & 6.0           & Added unit testing plan                                                                                                                                                                                                           \\
  2026-04-05          & 7.0           & Rev~1: Removed template comments; added unit test summary section; fixed typos and grammar; updated traceability matrices; addressed TA and peer feedback; added test environment details and timing requirements per peer review \\
  \bottomrule
\end{tabularx}

~\\

\newpage

\tableofcontents

\section{Symbols, Abbreviations, and Acronyms}

\renewcommand{\arraystretch}{1.2}
\begin{tabular}{l l}
  \toprule
  \textbf{symbol} & \textbf{description}                                              \\
  \midrule
  YOLO            & You Only Look Once -- an existing object detection model          \\
  CNN             & Convolutional Neural Network                                      \\
  Jetson          & abbreviation for specifically Jetson Nano Orin, which is the
  proposed board for this project
  \\
  STM32           & abbreviation for specifically STM32 microcontroller, which is the
  proposed microcontroller for this project                                           \\

  \bottomrule
\end{tabular}\\

\newpage

\pagenumbering{arabic}

This Verification and Validation (V\&V) plan outlines how the RoCam project
will demonstrate that the system is built correctly (verification) and fits
stakeholder needs (validation). It anchors testing to the SRS, architecture
(MG), and detailed design (MIS), ensuring every requirement is traceable where
possible. \textbf{Unit testing and system testing are performed;} only field
testing is out of scope. Verification relies on structured internal reviews and
checklists, plus unit tests and system tests, to confirm functional behavior
and safety assumptions. Validation is limited to a planned review with the
McMaster Rocketry Team (stakeholders) to collect feedback; no task-based
inspection, classmate review, or supervisor validation is conducted. Roles and
responsibilities are assigned to team members only, with feedback from the
rocketry team captured and tracked. Priorities emphasize correctness,
performance, reliability, and safety; lower-priority or out-of-scope items
(e.g., exhaustive usability studies, independent verification of third-party
components, and field testing) are noted to keep the plan feasible. The
document concludes with traceability tables, a unit test summary, and
appendices for parameters and optional surveys---forming a clear, practical
roadmap for the evidence of quality that can be gathered under these
constraints.

\section{General Information}

\subsection{Summary}

The RoCam system is a ground-based camera and gimbal assembly designed to
autonomously track model rockets during launch and flight. It utilizes a
high-resolution camera coupled with a motorized gimbal to maintain visual lock
on the rocket, providing real-time data and video feed. The system consists of
a STM32 based motion controller that interfaces with the physical gimbal; a
Jetson compute module that processes the camera feed; and a react frontend that
displays the video preview and allows the user to control the gimbal. The
primary purpose of RoCam is to replace the traditional unreliable manual
tracking methods with an automated solution to produce high-quality flight data
and video footage. To better assist model rocket teams across Canada to conduct
successful launches by providing an affordable and reliable system that can be
easily track rocket launch movements.

\subsection{Objectives}
The primary objective of this Verification and Validation V\&V Plan is to
establish confidence in the correctness, performance, and reliability of the
RoCam system by ensuring that all verification and validation activities are
thorough, traceable, and aligned with project requirements. The plan emphasizes
unit testing, system testing, and internal review to confirm that functional
and performance criteria are addressed where possible, with improvements
tracked through GitHub Issues. Validation focuses on a planned review with the
McMaster Rocketry Team (stakeholders) to collect feedback; no field testing or
supervisor validation is performed. These qualities—correctness, performance,
reliability, and traceability—are fundamental to the success of the V\&V
process, ensuring that the RoCam system is both technically sound and
operationally dependable. Certain objectives, such as independent verification
of third-party hardware and extended usability testing, are excluded due to
scope and resource limitations, with reliance placed on manufacturer and
stakeholder validation.

\subsection{In-Scope Objectives}
Verification and validation efforts will prioritize the following qualities:
\begin{itemize}
  \item \textbf{Correctness:} Ensure the reliability of tracking and control
        algorithms through testing and validation.
  \item \textbf{Performance:} Achieve real-time processing of 1080p 60\,fps
        video and maintain minimal gimbal control latency.
  \item \textbf{Safety:} Guarantee that the system operates safely under
        assumed conditions, with fail-safe mechanisms.
\end{itemize}

\subsection{Out-of-Scope Objectives}
The following objectives are excluded from this project due to time, budget,
and resource constraints:
\begin{itemize}
  \item \textbf{Field Testing:} Live rocket launches, hardware field conditions,
        and external reviewer validation at launch sites are out of scope.
  \item \textbf{Third-Party Component Verification:} External hardware such as the
        gimbal controller and camera are assumed to perform according to their
        manufacturer specifications.
  \item \textbf{Scalability Beyond Small-Scale Rockets:} Tracking a
        bullet or extremely difficult
        targets beyond a model rocket is outside the scope of this project.
\end{itemize}

By defining these priorities and exclusions, the RoCam project ensures that
available effort is directed toward validating the most critical system
qualities: \textbf{performance, accuracy, and operational safety}.

\subsection{Extras}

Beyond the core objectives, the RoCam project includes several extras that
enhance its quality and usability:
\begin{itemize}
  \item \textbf{STM32 control circuit design:} Development of a
        custom control circuit to interface the gimbal actuators with
        the host computer, enabling precise motion control and feedback sensing.
  \item \textbf{User Instructional Video:} Creation of a short instructional
        video demonstrating system setup, calibration, and operation to support
        new users and ensure safe deployment.
\end{itemize}

These extras aim to improve the overall usability, accessibility, and
documentation quality of the final deliverable, reinforcing RoCam's emphasis on
practical integration between computer vision software, mechanical control, and
end-user experience.

\subsection{Relevant Documentation}

\textbf{Software Requirements Specification (\SRS~\cite{SRS}):}
The SRS document serves as an outline for the VnV plan. The SRS document
contains specific functional and non-functional requirements, performance
criteria, and constraints. The SRS document also contains detailed stakeholder
roles and scope of the work. Aligning the VnV plan with the SRS document ensures
that all specified requirements are adequately verified and validated.

\textbf{Software Architecture MG (\MG~\cite{MG}) }
The MG document provides details of our system architecture. To effectively design
tests to validate and verify the system, it is important to understand the
architecture and functionality of each module.

\textbf{Software Detailed Design Document (MIS~\cite{MIS}):}
The MIS document provides detailed design of our system. It provides critical information
about the implementation details, including data structures, algorithms, and interfaces.
This document is essential for efficient test design to specific modules.

\section{Plan}

This section outlines the overall Verification and Validation (V\&V) plan for
the \textbf{RoCam} project. It defines the structure, responsibilities, and
methods used to verify that the system meets its functional and non-functional
requirements and to validate that it fulfills stakeholder needs. The section
covers the verification of the SRS, design, and implementation via internal
reviews, unit testing, and system testing, as well as the validation strategy
through a planned stakeholder review with the McMaster Rocketry Team (no field
testing). Its purpose is to provide a clear, organized framework ensuring that
all aspects of RoCam are systematically tested, reviewed, and proven reliable
before final deployment.

\subsection{Verification and Validation Team}

\begin{table}[H]
  \centering
  \caption{Verification and Validation Team Responsibilities}
  \label{tab:vnv_team}
  \begin{tabular}{|p{3.5cm}|p{3.5cm}|p{8cm}|}
    \hline
    \textbf{Name}         & \textbf{Role}              & \textbf{V\&V Responsibilities} \\ \hline

    \textbf{Zifan Si}     & Front End Interface Tester &
    Responsible for verifying and validating the usability and functionality of the
    react frontend interface.
    Ensure usability, robustness of frontend features.                                  \\ \hline

    \textbf{Jianqing Liu} & Team Lead                  &
    Conduct system validation. Validate the functionality and
    performance of the system against requirements and use cases. Validate system
    design and functionality.                                                           \\ \hline

    \textbf{Mike Chen}    & CV Tester                  &
    Verify and validate computer vision pipeline performance, including object detection
    and tracking accuracy.                                                              \\ \hline

    \textbf{Xiaotian Lou} & Test automation developer  &
    Implement automated testing frameworks in all software components. Implement
    API tests across different modules.
    \\ \hline

  \end{tabular}
\end{table}

\subsection{SRS Verification}

The SRS verification will focus on verifying whether the system design
satisfies all functional and non-functional requirements, stakeholder needs,
scope of work, and constraints specified. The SRS verification will be
conducted through the following steps:

\subsubsection{Functional and Non-Functional Requirement Review}
Since the design verification will evaluate in detail whether the system aligns
with the functional and non-functional requirements, the SRS verification will
focus on verifying the system design as a whole against the requirements. It
will be implemented as system tests in the final stages of development.

\subsubsection{Stakeholder Review}
Since our team leader is the control team lead of the McMaster Rocketry Team,
he serves as the primary reviewer of the system design and implementation and
acts as the communication bridge between the team and the stakeholders.

\subsubsection{Constraints and Scope of Work Review}
Several internal review sessions will be conducted to ensure that our system
design and implementation are within the defined scope of work and constraints
specified in the SRS document.

\begin{table}[H]
  \centering
  \caption{SRS Verification Checklist}
  \label{tab:srs_verification_checklist}
  \renewcommand{\arraystretch}{1.3}
  \setlength{\tabcolsep}{8pt}
  \begin{tabular}{|p{4.4cm}|p{10.1cm}|}
    \hline
    \textbf{Phase}                                  & \textbf{Verification Activities}                                                              \\ \hline

    Functional \& Non-Functional Requirement Review & $\square$ \; Review all
    requirements against current system design                                                                                                      \\[-2pt]
                                                    & $\square$ \; Create automated system tests for functional requirements                        \\[-2pt]
                                                    & $\square$ \; Measure performance, safety, and reliability against non-functional requirements \\[-2pt]
                                                    & $\square$ \; Record test results and address any failed cases                                 \\[-2pt]
                                                    & $\square$ \; Update traceability matrix after verification                                    \\ \hline

    Stakeholder Review                                                                                                                              \\[-2pt]
                                                    & $\square$ \; Present verified results and collect feedback                                    \\[-2pt]
                                                    & $\square$ \; Document required changes based on feedback                                      \\[-2pt]
                                                    & $\square$ \; Obtain stakeholder confirmation or approval                                      \\ \hline

    Constraints \& Scope of Work Review             & $\square$ \; Verify implementation stays
    within project scope                                                                                                                            \\[-2pt]
                                                    & $\square$ \; Check compliance with all technical and safety constraints                       \\[-2pt]
                                                    & $\square$ \; Review design against timeline and budget limits                                 \\[-2pt]
                                                    & $\square$ \; Update records if scope or constraints change                                    \\ \hline

    System Verification                             & $\square$ \; Conduct system tests covering all system
    components                                                                                                                                      \\[-2pt]
                                                    & $\square$ \; Validate full system behavior against SRS requirements                           \\[-2pt]
                                                    & $\square$ \; Log and resolve all system test issues                                           \\[-2pt]
                                                    & $\square$ \; Prepare final verification report for stakeholder review                         \\ \hline

    Safety Review                                   & $\square$ \; Review safety requirements against Hazard Analysis
    \\ \hline
  \end{tabular}
\end{table}

\subsection{Design Verification}

The design verification process for the \textbf{RoCam} project ensures that the
proposed architecture, algorithms, and hardware interfaces satisfy all
functional and non-functional requirements before implementation. Design
verification activities will be conducted through structured design reviews,
peer evaluations, and the use of detailed verification checklists. The goal is
to confirm that the design meets the intended performance, reliability, and
safety objectives defined in the SRS.

\subsubsection*{Verification Plan}
The design verification will proceed through internal review only (no classmate
or supervisor review, due to contact and resource limitations):
\begin{enumerate}
  \item \textbf{Internal Design Review:}
        Each subsystem (computer vision, control, and UI)
        will undergo an internal review by all four team members
        (Zifan Si, Jianqing Liu, Mike Chen, Xiaotian Lou).
        All members must agree on the designed architecture, algorithms,
        and interfaces. Interfaces must be clearly defined for scalability.
\end{enumerate}

\subsubsection*{Verification Artifacts}
Design verification will produce the following artifacts:
\begin{itemize}
  \item Verified design diagrams and interface specifications.
  \item Internal review notes and agreed-upon design decisions.
  \item Updated design documents reflecting corrections and improvements.
  \item Completed design verification checklists
        (Table~\ref{tab:design_verification_checklist}).
\end{itemize}

\subsubsection*{Design Verification Checklist}
Table~\ref{tab:design_verification_checklist} will be used during design reviews to ensure
completeness and consistency across all modules:

\newpage

\begin{table}[H]
  \centering
  \caption{Design Verification Checklist}
  \label{tab:design_verification_checklist}
  \renewcommand{\arraystretch}{1.3}
  \setlength{\tabcolsep}{8pt}
  \begin{tabular}{|p{4.4cm}|p{10.1cm}|}
    \hline
    \textbf{Phase}         & \textbf{Verification Activities}                                                     \\ \hline

    Internal Design Review & $\square$ \; Review subsystem architectures (vision,
    control, UI) with all team members                                                                            \\[-2pt]
                           & $\square$ \; Confirm APIs that are agreed upon                                       \\[-2pt]
                           & $\square$ \; Verify all module interfaces are clearly defined and documented in code \\[-2pt]
                           & $\square$ \; Ensure design supports scalability and modularity                       \\[-2pt]
                           & $\square$ \; Record review comments and update design diagrams                       \\ \hline

    Verification Artifacts & $\square$ \; Store verified design diagrams and
    updated interface specifications                                                                              \\[-2pt]
                           & $\square$ \; Archive internal review notes and design decisions                      \\[-2pt]
                           & $\square$ \; Maintain a design document in the project repository                    \\[-2pt]
                           & $\square$ \; Complete and file the final design verification checklist               \\ \hline
  \end{tabular}
\end{table}

\noindent
By following this structured verification plan and checklist, the team
ensures that the RoCam design is both technically sound and fully aligned
with its performance and safety goals prior to implementation.

\subsection{Verification and Validation Plan Verification}

The verification of the verification and validation (VnV) plan for the
\textbf{RoCam} project will be verified through internal review and unit
testing only. Each iteration of feedback and improvements will be tracked with
GitHub Issues. The focus of the verification plan is to ensure functionality
and performance metrics are adequately covered in design where possible.

For the validation plan, the only validation activity is a planned review with
the McMaster Rocketry Team (stakeholders) to collect feedback. No field tests
or supervisor validation are performed; unit testing and system testing are
conducted.

% Preamble:
% \usepackage{array}
% \usepackage{float}
% \usepackage{amssymb}

\begin{table}[H]
  \centering
  \caption{RoCam Verification and Validation (V\&V) Checklist}
  \label{tab:rocam_vnv_checklist}
  \renewcommand{\arraystretch}{1.3}
  \setlength{\tabcolsep}{8pt}
  \begin{tabular}{|p{4.4cm}|p{10.1cm}|}
    \hline
    \textbf{Phase}                 & \textbf{Actionable Verification and Validation Activities}                               \\ \hline

    Verification Plan Review       & $\square$ \; Conduct internal reviews of the V\&V
    plan with team                                                                                                            \\[-2pt]
                                   & $\square$ \; Ensure functionality and performance metrics are fully covered              \\[-2pt]
                                   & $\square$ \; Track feedback and revisions using GitHub Issues                            \\[-2pt]
                                   & $\square$ \; Confirm all review comments are resolved before approval                    \\ \hline

    Automated Verification Testing & $\square$ \; Develop automated tests for
    vision, control, and UI components                                                                                        \\[-2pt]
                                   & $\square$ \; Verify software behavior under simulated conditions                         \\[-2pt]
                                   & $\square$ \; Record test results and compare against performance criteria                \\[-2pt]
                                   & $\square$ \; Update test scripts based on identified defects or regressions              \\ \hline

    Continuous Monitoring          & $\square$ \; Review unit test outcomes after each
    iteration                                                                                                                 \\[-2pt]
                                   & $\square$ \; Maintain test evidence and validation summaries in repository               \\[-2pt]
                                   & $\square$ \; Update the V\&V plan to reflect system improvements and new unit test cases \\ \hline
  \end{tabular}
\end{table}

\subsection{Implementation Verification}

The implementation verification for the \textbf{RoCam} project ensures that the
developed software and integrated hardware components correctly realize the
verified design. This phase confirms that the implemented system satisfies the
functional and non-functional requirements defined in the SRS through a
combination of structured testing, static verification, and peer review.

\subsubsection*{Dynamic Verification}
Due to resource and contact limitations, \textbf{only field testing} is out of
scope. Unit testing and system testing are performed.

\begin{itemize}
  \item \textbf{Unit Testing:}
        Each software module (vision processing, gimbal control, data
        recording, and UI management) will undergo unit testing to verify
        functional correctness in isolation. This is the only form of
        dynamic testing conducted.
\end{itemize}

\noindent\textit{Out of scope:} Field testing is not performed due to the
constraints noted in this plan.

\subsubsection*{Static Verification}
Static verification will be performed throughout the implementation
phase to ensure that the codebase meets quality, safety,
and maintainability standards before execution.
Planned techniques include:

\begin{itemize}
  \item \textbf{Code Walkthroughs:}
        Internal code walkthroughs where each subsystem lead presents core
        implementation logic and design decisions to the team, enabling
        identification of potential design or coding issues.

  \item \textbf{Code Inspection:}
        Team members will review each other’s pull requests prior to
        merging into the main branch.
        Reviews will focus on coding standard compliance, potential
        logic errors, exception handling, and safety-critical sections that align
        with the Hazard Analysis
        (e.g., gimbal actuation and state transitions).

  \item \textbf{Static Analysis Tools:}
        Where applicable (e.g., on the frontend), static analysis may be used to
        detect common implementation defects. Results will be logged and reviewed
        as part of the verification record.

  \item \textbf{Documentation Review:}
        The team will verify that all inline comments, module-level
        documentation, and API references are consistent with the design
        specification and user documentation requirements.
\end{itemize}

\noindent
Together, these verification activities ensure that the RoCam
implementation is thoroughly tested, statically verified, and
validated against both its design and real-world performance expectations.

\subsection{Automated Testing and Verification Tools}

The RoCam project adopts a structured automated testing approach that
prioritizes maintainability, integration compatibility, and developer
productivity. The selection of testing frameworks was guided by three key
criteria: (1) \textbf{language and runtime fit}, ensuring each tool aligns with
its respective language and runtime; (2) \textbf{ease of automation}, to
integrate smoothly with CI/CD pipelines; and (3) \textbf{scalability}, so that
test coverage can expand as the project evolves. These principles ensure a
consistent and reliable testing environment across frontend and end-to-end
validation stages.

Each subsystem of RoCam employs unit testing and coverage tools appropriate to
its implementation language. The frontend (React/TypeScript) uses a unit
testing framework and coverage reporting suited to that stack. System-level and
Web UI testing, where in scope, use tooling that can coordinate frontend and
hardware layers. Specific frameworks are chosen to support CI/CD integration
and to allow test coverage to scale as the project evolves.

\textbf{Code Coverage and Reporting Plan:} \\
Coverage analysis is conducted using tools that provide line, branch, and
function-level metrics for each language stack, with results suitable for
unified tracking and CI integration.

\textbf{Linters and Formatters:} \\
To maintain coding standards and prevent stylistic or logical errors, RoCam enforces
consistent linting and formatting on the frontend:
\begin{itemize}
  \item \textbf{Frontend (TypeScript):} Uses \textbf{ESLint} and \textbf{Prettier} to
        maintain uniform style rules and detect syntax or logical inconsistencies.
\end{itemize}

\subsection{Software Validation}

Software validation confirms that \textbf{RoCam} satisfies stakeholder needs
and accurately implements the intended use cases (``building the right
system''). Validation complements verification and is referenced to the SRS
goals and stakeholder roles.

\subsubsection*{Stakeholder Review (McMaster Rocketry Team)}
Validation is limited to a planned review with the McMaster Rocketry Team
(stakeholders). The following is \textbf{in scope}:
\begin{itemize}
  \item \textbf{Stakeholder Review Session:} A scheduled review with the
        McMaster Rocketry Team to present the current system (or design/SRS),
        discuss scope, priorities, and acceptance criteria, and collect
        feedback. Revisions to the plan or implementation will be made based
        on their feedback.
\end{itemize}

\noindent\textbf{Out of scope} (not conducted): Task-based inspection by
stakeholders, classmate or peer review, supervisor validation, Rev~0/supervisor
demo, and formal user (operator) sessions. These are excluded due to contact
and resource limitations.

\subsubsection*{External Data and Bench Validation}
\begin{itemize}
  \item \textbf{Reference Footage:} Where live launches are unavailable,
        public model-rocket videos and internally recorded dry-runs are used
        to validate detection/tracking logic and reacquisition behavior.
  \item \textbf{Synthetic Scenarios:} Scripted clips with
        occlusion/smoke/backlight and varying apparent rocket sizes validate
        robustness claims prior to field tests.
\end{itemize}

\subsubsection*{Acceptance and Traceability}
\begin{itemize}
  \item \textbf{Acceptance Criteria:} Derived from the SRS (e.g.,
        sustained 1080p\,60\,fps I/O, end-to-end latency $\leq 120$\,ms,
        successful tracking through nominal flight phases, recording with
        angle overlays).
  \item \textbf{Traceability:} Validation tests map to SRS goals and
        stakeholder needs; a table links each validation activity to the
        corresponding SRS item(s) and planned evidence (videos, logs, checklists).
\end{itemize}

\noindent
If no suitable external datasets are available for certain scenarios,
this limitation will be stated explicitly; validation in this plan relies
on the planned stakeholder review with the McMaster Rocketry Team and, where
applicable, unit testing and internal review only. This section also references the SRS verification section for
consistency checks between validated behaviors and specified requirements.

\section{System Tests}

\textbf{Scope note:} Due to resource and contact limitations, \textbf{only field
  testing} is out of scope. Unit tests and system tests
described in this section are executed. Evidence of correctness comes from unit
tests, system tests, and the planned stakeholder review
with the McMaster Rocketry Team.

\subsection{Tests for Functional Requirements}

The following manual tests are designed to evaluate the robustness of the
Tracking Camera System. Each test includes its initial state, input, expected
output, justification (test case derivation), and procedure. Tests are
organized by functional area. The intent is to verify compliance with the
functional requirements (FR) while deliberately stressing the system to expose
edge cases.

\subsubsection{Area of Testing: Video Acquisition and Output}

\subsubsection*{\tid{test-acquire-frame}: Camera Frames are Acquired at Startup}
\begin{itemize}
  \item \textbf{Control:} Manual
  \item \textbf{Initial State:} System powered off; camera connected; HDMI monitor attached.
  \item \textbf{Input:} Power on the system and allow it to fully initialize.
  \item \textbf{Expected Output:} Within 10s of boot, logs show that the camera is acquiring frames.
  \item \textbf{Test Case Derivation:} The system's operation depends on a valid video feed.
  \item \textbf{How Test Will Be Performed:} Boot the device; observe the program logs.
\end{itemize}

\subsubsection*{\tid{test-disconnect-camera}: Camera Disconnect During Operation}
\begin{itemize}
  \item \textbf{Control:} Manual
  \item \textbf{Initial State:} System initialized; preview and HDMI output active.
  \item \textbf{Input:} Physically disconnect the camera from the system.
  \item \textbf{Expected Output:} On disconnect, the UI and HDMI output show a clear error message; no crash.
  \item \textbf{Test Case Derivation:} When the camera is disconnected, the system should gracefully handle the loss without crashing. And it should notify the operator that the camera is disconnected.
  \item \textbf{How Test Will Be Performed:} Disconnect the camera from the system; observe the UI and HDMI output for error messages.
\end{itemize}

\subsubsection*{\tid{test-hdmi-output}: HDMI Output with Info Overlay}
\begin{itemize}
  \item \textbf{Control:} Manual
  \item \textbf{Initial State:} System initialized; idle state; preview and HDMI output active.
  \item \textbf{Input:} Move the gimbal manually through the UI.
  \item \textbf{Expected Output:} HDMI shows 1080p60 camera feed with a text overlay of current pan/tilt angles.
  \item \textbf{Test Case Derivation:} The system should be able to display the current camera feed and the pan/tilt angles for the viewers
  \item \textbf{How Test Will Be Performed:} Move the gimbal manually, and visually verify overlay values change smoothly and match commanded motion
\end{itemize}

\subsubsection*{\tid{test-disconnect-hdmi}: HDMI Disconnect}
\begin{itemize}
  \item \textbf{Control:} Manual
  \item \textbf{Initial State:} System initialized; preview and HDMI output active.
  \item \textbf{Input:} Physically disconnect the HDMI cable.
  \item \textbf{Expected Output:} On disconnect, the UI shows a clear error message; no crash.
  \item \textbf{Test Case Derivation:} When the HDMI cable is disconnected, the system should gracefully handle the loss without crashing. And it should notify the operator that the HDMI cable is disconnected.
  \item \textbf{How Test Will Be Performed:} Disconnect the HDMI cable from the system; observe the UI for error messages.
\end{itemize}

\subsubsection*{\tid{test-preview-real-time}: Preview with Real-Time Info}
\begin{itemize}
  \item \textbf{Control:} Manual
  \item \textbf{Initial State:} System initialized; idle state; preview and HDMI output active.
  \item \textbf{Input:} Move the gimbal manually through the UI.
  \item \textbf{Expected Output:} Preview in the UI updates fluidly (target $\geq$15fps) and displays current pan/tilt angles.
  \item \textbf{Test Case Derivation:} The system should allow the camera operator to see the current pan/tilt angles and HDMI output in real-time.
  \item \textbf{How Test Will Be Performed:} Move the gimbal manually, and visually verify the preview in the UI updates fluidly and displays the current pan/tilt angles.
\end{itemize}

\subsubsection*{\tid{test-network-disconnect}: Network Disconnect While Viewing UI}
\begin{itemize}
  \item \textbf{Control:} Manual
  \item \textbf{Initial State:} System initialized; preview and HDMI output active.
  \item \textbf{Input:} Disconnect the system from the network for 30s; reconnect.
  \item \textbf{Expected Output:} The system keeps functioning as normal without network; UI shows a clear error message; upon reconnection, the UI session restores automatically within 5s without reboot; no data loss or crash.
  \item \textbf{Test Case Derivation:} The system needs to be resilient against transient network failures in field conditions.
  \item \textbf{How Test Will Be Performed:} Manually disconnect and connect the system from the network; observe the UI for error messages.
\end{itemize}

\subsubsection{Area of Testing: Recording}

\subsubsection*{\tid{test-recording}: Start and Stop Recording}
\begin{itemize}
  \item \textbf{Control:} Manual
  \item \textbf{Initial State:} System initialized; idle state; sufficient disk space.
  \item \textbf{Input:} Press "Start Recording", wait 10s, then press "Stop Recording".
  \item \textbf{Expected Output:} A 10 second 1080p60 video file is created and saved to the disk.
  \item \textbf{Test Case Derivation:} The system should be able to start and stop recording.
  \item \textbf{How Test Will Be Performed:} Press "Start Recording", wait 10s, then press "Stop Recording".
\end{itemize}

\subsubsection*{\tid{test-disk-full}: Disk Full Behavior During Recording}
\begin{itemize}
  \item \textbf{Control:} Manual
  \item \textbf{Initial State:} System initialized; idle state; storage nearly full (simulate by filling disk leaving $<1$min recording capacity).
  \item \textbf{Input:} Start recording and allow disk to fill.
  \item \textbf{Expected Output:} UI shows a clear "Storage Full" alert \textbf{within 2 seconds} of exhaustion; current recording is stopped and the file is saved to the disk; no crash.
  \item \textbf{Test Case Derivation:} The system should be able to handle storage exhaustion gracefully.
  \item \textbf{How Test Will Be Performed:} Fill the disk with data until it is nearly full; start recording; observe the UI for the "Storage Full" alert; check the disk for the saved file.
\end{itemize}

\subsubsection*{\tid{test-manage-recordings}: List and Download Recordings}
\begin{itemize}
  \item \textbf{Control:} Manual
  \item \textbf{Initial State:} At least two prior recordings exist.
  \item \textbf{Input:} Open recordings list; download the latest file and an older file; delete a small test clip.
  \item \textbf{Expected Output:} List shows correct entries (name, timestamp, duration); downloads succeed; delete removes the selected item only.
  \item \textbf{Test Case Derivation:} The system should be able to list, delete, and download recordings.
  \item \textbf{How Test Will Be Performed:} Open the recordings list; download the latest file and an older file; delete a small test clip; check the list for the correct entries; check the disk for the saved files.
\end{itemize}

\subsubsection{Area of Testing: Tracking}

\subsubsection*{\tid{test-manual-control}: Manual Gimbal Control in Idle}
\begin{itemize}
  \item \textbf{Control:} Manual
  \item \textbf{Initial State:} System initialized; idle state; preview and HDMI output active.
  \item \textbf{Input:} Issue manual pan and tilt commands via UI.
  \item \textbf{Expected Output:} The gimbal follows commands smoothly.
  \item \textbf{Test Case Derivation:} The operator should be able to manually control the gimbal in idle state.
  \item \textbf{How Test Will Be Performed:} Issue manual pan and tilt commands via UI; observe the gimbal movement.
\end{itemize}

\subsubsection*{\tid{test-transition-to-tracking}: Transition to Tracking When Rocket Detected}
\begin{itemize}
  \item \textbf{Control:} Manual
  \item \textbf{Initial State:} System initialized; idle state; preview and HDMI output active.
  \item \textbf{Input:} Arm the system via UI, and waving a model rocket in front of the camera.
  \item \textbf{Expected Output:} The system transitions to tracking state as soon as the rocket is detected; the UI shows the system is in tracking state
  \item \textbf{Test Case Derivation:} The system should be able to transition to tracking state when a rocket is detected.
  \item \textbf{How Test Will Be Performed:} Arm the system via UI, and waving a model rocket in front of the camera; observe the UI for the system to transition to tracking state.
\end{itemize}

\subsubsection*{\tid{test-tracking}: Keep Rocket in Frame While Tracking}
\begin{itemize}
  \item \textbf{Control:} Manual
  \item \textbf{Initial State:} System initialized; tracking state; preview and HDMI output active.
  \item \textbf{Input:} Wave a model rocket (approx.\ 30\,cm length) at 1--3\,m distance from the camera in an indoor lab environment with standard overhead lighting. Move the target at moderate speed (approx.\ 0.5--1\,m/s lateral) across the camera's field of view.
  \item \textbf{Expected Output:} Gimbal should continuously adjust its position to keep the rocket within the center 80\% of the frame for the duration of the test (minimum 10\,s of sustained tracking).
  \item \textbf{Test Case Derivation:} The system should be able to keep the rocket in the frame while tracking. This test verifies FR-7.
  \item \textbf{How Test Will Be Performed:} In a controlled indoor environment, wave the model rocket at the specified distance and speed in front of the camera; observe the gimbal movement tracks the target; observe the preview in the UI to confirm the rocket stays in frame.
\end{itemize}

\subsubsection*{\tid{test-return-to-idle}: Return to Idle When Rocket is Lost}
\begin{itemize}
  \item \textbf{Control:} Manual
  \item \textbf{Initial State:} System initialized; tracking state; preview and HDMI output active.
  \item \textbf{Input:} Occlude the model rocket.
  \item \textbf{Expected Output:} Within 2s of sustained loss, system transitions to idle state; ceases autonomous gimbal commands; UI indicates "target lost".
  \item \textbf{Test Case Derivation:} The system should be able to return to idle state when the rocket is lost.
  \item \textbf{How Test Will Be Performed:} Occlude the model rocket; observe the UI for the system to transition to idle state.
\end{itemize}

\subsubsection*{\tid{test-exit-gimbal-range}: Rocket Exits Gimbal Range}
\begin{itemize}
  \item \textbf{Control:} Manual
  \item \textbf{Initial State:} System initialized; tracking state; preview and HDMI output active.
  \item \textbf{Input:} Move the model rocket outside of the gimbal's range of motion.
  \item \textbf{Expected Output:} The gimbal moves until it reaches its mechanical limit, then stops moving further; no errors or warnings are presented to the user; system remains responsive.
  \item \textbf{Test Case Derivation:} The system must safely handle objects that leave the gimbal's range without faulting or reporting spurious errors.
  \item \textbf{How Test Will Be Performed:} Move the model rocket outside of the gimbal's range of motion; observe the gimbal reaches its mechanical end and stops; verify that no error messages appear in the UI and that all other system functions remain operational.
\end{itemize}

\subsubsection{Field Testing (Out of Scope)}

Field testing is \textbf{out of scope} for this project due to resource and
contact limitations. The following is retained for reference only.

If field testing were to be conducted, it would be done by the McMaster
Rocketry Team (client), who would be provided with the manual and would operate
the system at a model rocket launch event. The development team would be
present to assist. After the event, the client could fill out a survey
(Appendix 6.2). Under the current plan, no field testing is performed; unit
testing, system testing, and a planned stakeholder review with the rocketry
team are conducted.

\subsection{Tests for Nonfunctional Requirements}

\subsubsection{UI Walkthrough}

The UI walkthrough will be conducted by the capstone team members themselves
before the field test.

\subsubsection*{\tid{test-consistent-units}: Consistent and Changeable Units}
\begin{itemize}
  \item \textbf{Type:} Manual
  \item \textbf{Initial State:} System initialized; idle state; preview and HDMI output active.
  \item \textbf{Input/Condition:} Go through the UI twice, once with metric units, once with imperial units.
  \item \textbf{Output/Result:} The units are consistent and changeable.
  \item \textbf{How Test Will Be Performed:} The capstone team members will go through the UI twice, once with metric units, once with imperial units; observe the units in the UI.
\end{itemize}

\subsubsection*{\tid{test-bilingual-ui}: Bilingual UI}
\begin{itemize}
  \item \textbf{Type:} Manual
  \item \textbf{Initial State:} System initialized; language = English.
  \item \textbf{Input/Condition:} Switch UI language to French and back; navigate all UI sections, dialogs, and toasts.
  \item \textbf{Output/Result:} All visible strings/localized assets appear in the selected language; no truncation/overflow; no mixed-language strings; hotkeys remain functional; date/time/number formats localize correctly.
  \item \textbf{How Test Will Be Performed:} Perform a complete UI walkthrough in French; capture screenshots of each page/dialog; switch back to English and repeat spot checks.
\end{itemize}

\subsubsection*{\tid{test-immediate-ui-feedback}: Immediate UI Feedback}
\begin{itemize}
  \item \textbf{Type:} Manual
  \item \textbf{Initial State:} System idle; UI ready.
  \item \textbf{Input/Condition:} Trigger representative actions: button press, toggle, menu selection, starting/stopping recording, arming/disarming, opening settings.
  \item \textbf{Output/Result:} A visible confirmation (pressed state, spinner, toast, label change) appears within 0.2s for each action.
  \item \textbf{How Test Will Be Performed:} Observe the UI while performing each action
\end{itemize}

\subsubsection*{\tid{test-visual-indicators}: Rely on Visual Indicators}
\begin{itemize}
  \item \textbf{Type:} Manual
  \item \textbf{Initial State:} System initialized; speakers muted or disconnected.
  \item \textbf{Input/Condition:} Trigger statuses/alerts: armed, tracking, target lost, storage low/full, network disconnect, device error.
  \item \textbf{Output/Result:} Every status/alert is presented visually (icons, colors, banners, toasts) without relying on sound; no action-critical info is audio-only.
  \item \textbf{How Test Will Be Performed:} With audio disabled, induce each state/fault; verify a visible indicator appears and remains long enough to be noticed.
\end{itemize}

\subsubsection*{\tid{test-confirmation-prompts}: Confirmation Prompts for Safety-Critical Actions}
\begin{itemize}
  \item \textbf{Type:} Manual
  \item \textbf{Initial State:} System idle; gimbal powered.
  \item \textbf{Input/Condition:} Attempt potentially hazardous actions: force home while tracking, hard stop, high-speed slew, firmware update during armed state, factory reset.
  \item \textbf{Output/Result:} A blocking confirmation dialog appears for each action; dialog text states the specific risk and outcome; default action is “Cancel”; action proceeds only after explicit confirmation.
  \item \textbf{How Test Will Be Performed:} Trigger each action; capture screenshots of the prompt; verify risk wording is specific (not generic); verify no motion/change occurs until “Confirm” is pressed.
\end{itemize}

\subsubsection*{\tid{test-manual-idle-mode}: Manual Idle Mode / Override}
\begin{itemize}
  \item \textbf{Type:} Manual
  \item \textbf{Initial State:} System armed or tracking.
  \item \textbf{Input/Condition:} Press “Idle” (or equivalent) control while tracking.
  \item \textbf{Output/Result:} System transitions to idle within 1\,s; autonomous commands cease; manual controls become available; status clearly shows “Idle.”
  \item \textbf{How Test Will Be Performed:} Enter tracking, then invoke manual idle; observe immediate cessation of gimbal motion; verify state banner change and manual jog controls responsiveness.
\end{itemize}

\subsubsection*{\tid{test-no-collection-pii}: No Collection of PII}
\begin{itemize}
  \item \textbf{Type:} Manual
  \item \textbf{Initial State:} System initialized; default configuration.
  \item \textbf{Input/Condition:} Inspect all settings, logs, exported files, filenames/paths, and network-configurable fields.
  \item \textbf{Output/Result:} No fields request or store personal identifiers (names, emails, phone numbers, precise home addresses, faces with identity tags); logs contain only technical data; telemetry excludes PII.
  \item \textbf{How Test Will Be Performed:} Open each settings and log panel; download logs/recordings; inspect metadata (EXIF/JSON/CSV) for PII keys; verify any optional labels are device/session IDs, not personal data.
\end{itemize}

\subsubsection*{\tid{test-color-vision-deficiency}: Color-Vision Deficiency Accessibility}
\begin{itemize}
  \item \textbf{Type:} Manual
  \item \textbf{Initial State:} System initialized; all status indicators visible.
  \item \textbf{Input/Condition:} Review all status/alert states (normal, armed, tracking, warning, error) and interactive controls (enabled/disabled) under simulated Protanopia/Deuteranopia/Tritanopia or printed CVD reference swatches.
  \item \textbf{Output/Result:} Each state is distinguishable without relying on hue alone (uses shape, text labels, patterns, or contrast); contrast ratio \(\geq\) 4.5:1 for text/icons against background.
  \item \textbf{How Test Will Be Performed:} Cycle UI through all states; verify distinct icon shapes/labels; optionally view via a CVD simulator; record pass/fail with screenshots; confirm tooltips/textual labels exist for each color state.
\end{itemize}

\subsubsection{Code and Documentation Walkthrough}

\subsubsection*{\tid{test-manual-arm-mode}: Manual Arm Mode Only}
\begin{itemize}
  \item \textbf{Type:} Static
  \item \textbf{Initial State:} N/A
  \item \textbf{Input/Condition:} Inspect all state transition code and configuration defaults.
  \item \textbf{Output/Result:} The system only enters armed mode when explicitly commanded by the user, never automatically.
  \item \textbf{How Test Will Be Performed:} Review code for all references to the armed state, verify guards require user input.
\end{itemize}

\subsubsection*{\tid{test-gimbal-interface}: Gimbal Interface Documentation Review}
\begin{itemize}
  \item \textbf{Type:} Static
  \item \textbf{Initial State:} N/A
  \item \textbf{Input/Condition:} Review the interface specification with the client, including command set, telemetry, timing, and safety notes.
  \item \textbf{Output/Result:} The client confirms the interface documentation is complete, clear, and suitable for supporting multiple gimbal systems.
  \item \textbf{How Test Will Be Performed:} Conduct a walkthrough meeting with the client, check all required topics against a prepared checklist.
\end{itemize}

\subsubsection{Error Handling and Logging}

\subsubsection*{\tid{test-report-errors}: Report All Errors to the User}
\begin{itemize}
  \item \textbf{Type:} Manual
  \item \textbf{Initial State:} System initialized and idle.
  \item \textbf{Input/Condition:} Manually trigger hardware, communication, and software faults.
  \item \textbf{Output/Result:} Each error generates a visible and descriptive message on the user interface \textbf{within 2 seconds} of fault occurrence.
  \item \textbf{How Test Will Be Performed:} Disconnect devices, kill processes, and simulate faults while observing the UI for clear error reporting. Time the interval between fault injection and the appearance of the error notification; verify it does not exceed 2\,s.
\end{itemize}

\subsubsection*{\tid{test-stop-immediately}: Stop Immediately on Unrecoverable Error}
\begin{itemize}
  \item \textbf{Type:} Manual
  \item \textbf{Initial State:} System operating in tracking mode.
  \item \textbf{Input/Condition:} Introduce a fatal condition such as gimbal communication loss or corrupted control process.
  \item \textbf{Output/Result:} The system stops all automated operations and transitions to a safe idle state \textbf{within 0.5 seconds} of the unrecoverable fault.
  \item \textbf{How Test Will Be Performed:} Disconnect the gimbal or crash the tracking process; time the interval until motion ceases and the system reaches idle state; verify that the interval does not exceed 0.5\,s and that a clear unrecoverable error message is displayed.
\end{itemize}

\subsubsection*{\tid{test-validate-commands}: Validate All Commands}
\begin{itemize}
  \item \textbf{Type:} Manual
  \item \textbf{Initial State:} System idle.
  \item \textbf{Input/Condition:} Enter invalid or out-of-range user commands through the UI or API.
  \item \textbf{Output/Result:} The system rejects invalid commands and displays a validation error without sending them to the gimbal.
  \item \textbf{How Test Will Be Performed:} Attempt to send commands beyond gimbal limits or in unsupported formats and confirm that they are blocked with clear feedback.
\end{itemize}

\subsubsection*{\tid{test-log-operations}: Log All Operations}
\begin{itemize}
  \item \textbf{Type:} Manual
  \item \textbf{Initial State:} System initialized and user logged in.
  \item \textbf{Input/Condition:} Perform a series of user operations such as arming, tracking, and recording.
  \item \textbf{Output/Result:} Each action is logged with a timestamp and operation details in the system log.
  \item \textbf{How Test Will Be Performed:} After the field test, review the log file to confirm all actions and timestamps are recorded accurately.
\end{itemize}

\subsubsection{Field Testing (Out of Scope)}

Field testing and user surveys are out of scope; unit testing, system testing,
and the stakeholder review are performed.

\subsubsection*{\tid{test-user-survey}: User Survey (Out of Scope)}
\begin{itemize}
  \item \textbf{Type:} Static (out of scope)
  \item \textbf{Initial State:} N/A
  \item \textbf{Input/Condition:} N/A
  \item \textbf{Output/Result:} Completed user survey.
  \item \textbf{How Test Will Be Performed:} Not performed. If field testing were in scope, the client would fill out the survey after the field test. Under the current plan, unit testing, system testing, and stakeholder review are conducted.
\end{itemize}

\subsection{Traceability Between Test Cases and Requirements}

\begin{tabularx}{\textwidth}{p{5cm}X}
  \toprule {\bf Test ID}             & {\bf Requirements}       \\
  \midrule
  \tid{test-acquire-frame}           & FR-1, SLR-2              \\
  \tid{test-disconnect-camera}       & FR-1                     \\
  \tid{test-hdmi-output}             & FR-8, FR-9, INT-1, SLR-3 \\
  \tid{test-disconnect-hdmi}         & FR-8                     \\
  \tid{test-preview-real-time}       & FR-10                    \\
  \tid{test-network-disconnect}      & FR-10                    \\
  \tid{test-recording}               & FR-11                    \\
  \tid{test-disk-full}               & FR-11                    \\
  \tid{test-manage-recordings}       & FR-12                    \\
  \tid{test-manual-control}          & FR-2, FR-6               \\
  \tid{test-transition-to-tracking}  & FR-3, FR-4               \\
  \tid{test-tracking}                & FR-2, FR-3, FR-7         \\
  \tid{test-return-to-idle}          & FR-5                     \\
  \tid{test-exit-gimbal-range}       & FR-2, FR-7               \\
  \tid{test-consistent-units}        & EZ-1, PI-1               \\
  \tid{test-bilingual-ui}            & PI-2                     \\
  \tid{test-immediate-ui-feedback}   & EZ-2                     \\
  \tid{test-visual-indicators}       & EZ-3                     \\
  \tid{test-confirmation-prompts}    & EZ-4, UPR-3              \\
  \tid{test-manual-idle-mode}        & SCR-1                    \\
  \tid{test-no-collection-pii}       & LR-1                     \\
  \tid{test-color-vision-deficiency} & AR-1                     \\
  \tid{test-manual-arm-mode}         & SCR-4                    \\
  \tid{test-gimbal-interface}        & INT-2                    \\
  \tid{test-report-errors}           & RFR-1                    \\
  \tid{test-stop-immediately}        & RFR-2                    \\
  \tid{test-validate-commands}       & IR-1                     \\
  \tid{test-log-operations}          & AUR-1                    \\
  \bottomrule
\end{tabularx}

\section{Unit Test Description}

\subsection{Unit Test Summary}

This subsection provides an overview of the unit testing approach, frameworks,
module coverage, and expected coverage targets for the RoCam project.

\subsubsection*{Testing Approach}
Unit testing follows a bottom-up strategy: each leaf module identified in the
Module Guide~(\MG~\cite{MG}) is tested in isolation before integration. Tests
are designed to verify functional correctness, boundary conditions, error
handling, and key nonfunctional attributes (performance, robustness, thread
safety). External dependencies (hardware, third-party SDKs) are replaced by
mocks or stubs so that tests are deterministic and repeatable.

\subsubsection*{Frameworks and Tooling}
\begin{itemize}
  \item \textbf{Backend (Python) -- pytest:}
        All backend modules running on the Jetson are tested with
        \texttt{pytest}. Fixtures provide reusable mocked dependencies
        (gimbal, database, IPC connections). Coverage is measured with
        \texttt{pytest-cov}, producing line, statement, and branch metrics.
  \item \textbf{Frontend (React/TypeScript) -- Vitest:}
        All frontend UI modules are tested with \texttt{Vitest} (a
        Vite-native test runner compatible with the Jest API). Component
        tests use \texttt{@testing-library/react} for DOM assertions.
        Coverage is measured with Vitest's built-in V8 provider, producing
        line, statement, branch, and function metrics.
  \item \textbf{CI Integration:}
        Both test suites run automatically on every pull request targeting
        \texttt{main} via GitHub Actions. All tests must pass before a pull
        request can be merged.
\end{itemize}

\subsubsection*{Modules Covered}
\begin{itemize}
  \item \textbf{Backend (18 modules):}
        State Management, API Gateway, Recording Database, Gimbal
        Abstraction, Tracking, CV Process Management, Live Video Process
        Management, Transcode API, Transcode Process, CV Process, Common
        IPC, Common IPC Buffer, Common Watchdog, System Status, Common
        Utils, Recording-Only State Management, Main Entry Point, and SSE.
  \item \textbf{Frontend (10 module groups):}
        Preview, Manual Control, Recording Management, Configuration, API
        Client, Rocam Provider, Utility Formatters, Internationalization,
        App Routing, and individual UI Components (Navbar, SystemStatusCard,
        ControlsCard).
\end{itemize}

\subsubsection*{Out-of-Scope Modules}
The following modules are excluded from unit testing because they wrap
third-party or hardware-specific libraries:
\begin{itemize}
  \item Video Stream Abstraction (Nvidia Deepstream SDK)
  \item Serial Abstraction (pyserial)
  \item System Status hardware queries (jetson-stats)
\end{itemize}

\subsubsection*{Coverage Targets}
The project targets a minimum of \textbf{80\% line coverage} for both
backend and frontend codebases. Achieved coverage at the time of the VnV
Report is:
\begin{itemize}
  \item Backend: 88\% line coverage, 81.4\% branch coverage (313 tests passed)
  \item Frontend: 85.57\% line coverage, 70.35\% branch coverage (168 tests passed)
\end{itemize}
Branch coverage is lower because some defensive paths (e.g., hardware
error branches, rare race conditions) are difficult to trigger in a
mocked environment. These paths are partially addressed by system tests
and manual walkthroughs.

\subsection{Unit Testing Scope}

The RoCam system unit testing covers all leaf modules implemented by the
project team. The following modules are within the scope of unit testing:

\begin{itemize}
  \item \textbf{Backend (Jetson) Modules:}
        \begin{itemize}
          \item \mref{mState} (State Management Module)
          \item \mref{mAPI} (API Gateway Module)
          \item \mref{mRecord} (Recording Database Module)
          \item \mref{mGimbal} (Gimbal Abstraction Module)
          \item \mref{mTrack} (Tracking Module)
          \item \mref{mCVManagement} (CV Process Management Module)
          \item \mref{mVideoManagement} (Live Video Process Management Module)
          \item \mref{mTranscodeManagement} (Transcode Process Management Module)
          \item \mref{mCV} (CV Process Module)
          \item \mref{mVideo} (Live Video Process Module)
          \item \mref{mTranscode} (Transcode Process Module)
          \item Common utilities (IPC, IPC Buffer, Watchdog, System Status, Utils)
        \end{itemize}

  \item \textbf{Frontend (UI) Modules:}
        \begin{itemize}
          \item \mref{mPreview} (Preview Module)
          \item \mref{mManual} (Manual Control Module)
          \item \mref{mRecordMng} (Recording Management Module)
          \item \mref{mConfig} (Configuration Module)
          \item API client and network layer
          \item Utility formatters and internationalization
        \end{itemize}
\end{itemize}

\textbf{Out of Scope:}
\begin{itemize}
  \item \mref{mStream} (Video Stream Abstraction Module) --- Implemented by Nvidia Deepstream SDK (third-party)
  \item \mref{mSerial} (Serial Abstraction Module) --- Implemented by pyserial library (third-party)
  \item \mref{mStatus} (System Status Module) --- Hardware-dependent library (jetson-stats)
  \item Hardware-specific system tests (require actual Jetson hardware)
\end{itemize}

The unit testing strategy employs pytest for Python backend modules and Vitest
for TypeScript frontend modules. Tests are designed to verify functional
correctness, edge case handling, and error conditions for each module in
isolation, with external dependencies mocked appropriately.

\subsection{Tests for Functional Requirements}

Unit tests for functional requirements verify that each module correctly
implements its specified behavior. Tests are organized by module, with
references to the MIS and MG documents where applicable. All tests are
automated and implemented using pytest (Python backend) or Vitest (TypeScript
frontend).

\subsubsection{Module: State Management (\mref{mState})}

The State Management module handles system state transitions, manual gimbal
control, recording operations, and status aggregation. Tests verify state
machine behavior, command validation, and data management.

\begin{enumerate}

  \item{\tid{test-state-management-arm-disarm}\\}

        Type: Functional, Dynamic, Automatic

        Initial State: StateManagement instance with mocked dependencies

        Input: Call arm() and disarm() methods

        Output: Internal armed flag set appropriately; gimbal LED commands issued

        Test Case Derivation: Verifies FR-4 (arming/disarming functionality) and that
        state transitions trigger correct hardware responses

        How test will be performed: Automated via
        test\_state\_management\_class.py::TestArmDisarm

  \item{\tid{test-state-management-manual-move}\\}

        Type: Functional, Dynamic, Automatic

        Initial State: Disarmed system with gimbal at known position

        Input: manual\_move() with directions up/down/left/right

        Output: Gimbal set\_deg() called with correctly adjusted angles

        Test Case Derivation: Verifies FR-2 (manual control) and FR-6 (manual control
        only when disarmed)

        How test will be performed: Automated via
        test\_state\_management\_class.py::TestManualMove

  \item{\tid{test-state-management-manual-move-clamped}\\}

        Type: Functional, Dynamic, Automatic

        Initial State: Disarmed system with gimbal near limits

        Input: manual\_move() commands that would exceed gimbal limits

        Output: Angles clamped to valid ranges (tilt 0-90°, pan -45° to 45°)

        Test Case Derivation: Verifies IR-1 (command validation and clamping)

        How test will be performed: Automated via
        test\_state\_management\_class.py::TestManualMove

  \item{\tid{test-state-management-recording}\\}

        Type: Functional, Dynamic, Automatic

        Initial State: System with mocked database and CV process

        Input: start\_recording() and stop\_recording() calls

        Output: Recording allocated in database; CV process notified; recording ID
        tracked

        Test Case Derivation: Verifies FR-11 (recording start/stop functionality)

        How test will be performed: Automated via
        test\_state\_management\_class.py::TestRecording

  \item{\tid{test-state-management-download-pause}\\}

        Type: Functional, Dynamic, Automatic

        Initial State: Active recording session

        Input: on\_download\_start() and on\_download\_end() calls during recording

        Output: CV pipeline paused on first download, resumed after last download ends

        Test Case Derivation: Verifies correct pipeline management during recording
        downloads

        How test will be performed: Automated via
        test\_state\_management\_class.py::TestDownloadHandlers

  \item{\tid{test-bounding-box-collection}\\}

        Type: Functional, Dynamic, Automatic

        Initial State: Empty BoundingBoxCollection

        Input: Sequential CVData with bounding boxes and timestamps

        Output: Correct bbox retrieval within 40ms window; latest valid bbox returned;
        max 10 entries enforced

        Test Case Derivation: Verifies data structure behavior for tracking data
        management

        How test will be performed: Automated via test\_state\_management.py

\end{enumerate}

\subsubsection{Module: API Gateway (\mref{mAPI})}

The API Gateway module provides REST endpoints for UI communication. Tests
verify endpoint behavior, request validation, and proper delegation to
StateManagement.

\begin{enumerate}

  \item{\tid{test-api-endpoints-exist}\\}

        Type: Functional, Dynamic, Automatic

        Initial State: Flask test client with mocked StateManagement

        Input: HTTP requests to /api/arm, /api/disarm, /api/manual\_move,
        /api/manual\_move\_to, /api/recordings/start, /api/recordings/stop

        Output: HTTP 200 responses; correct methods called on StateManagement

        Test Case Derivation: Verifies FR-6, FR-10, FR-11 (API endpoints for control
        and recording)

        How test will be performed: Automated via test\_api.py and
        test\_api\_gateway.py

  \item{\tid{test-api-recordings-crud}\\}

        Type: Functional, Dynamic, Automatic

        Initial State: Flask test client with mocked database

        Input: GET /api/recordings, PATCH /api/recordings/<id>, DELETE
        /api/recordings/<id>

        Output: Recording list returned; rename and delete operations delegated to
        database

        Test Case Derivation: Verifies FR-12 (recording management)

        How test will be performed: Automated via test\_api.py::TestRecordingsList,
        TestRecordingsRename, TestRecordingsDelete

  \item{\tid{test-api-validation}\\}

        Type: Functional, Dynamic, Automatic

        Initial State: Flask test client

        Input: PATCH /api/recordings/<id> with missing or invalid new\_name

        Output: HTTP 400 Bad Request with appropriate error message

        Test Case Derivation: Verifies IR-1 (input validation)

        How test will be performed: Automated via test\_api.py::TestRecordingsRename

\end{enumerate}

\subsubsection{Module: Recording Database (\mref{mRecord})}

The Recording Database module manages recording metadata, file storage, and
retrieval operations.

\begin{enumerate}

  \item{\tid{test-database-allocate}\\}

        Type: Functional, Dynamic, Automatic

        Initial State: Empty temporary directory

        Input: allocate\_recording() call

        Output: Recording directory created with meta.json; RecordingInfo returned with
        valid paths

        Test Case Derivation: Verifies FR-11 (recording allocation)

        How test will be performed: Automated via
        test\_database.py::TestAllocateRecording

  \item{\tid{test-database-list-recordings}\\}

        Type: Functional, Dynamic, Automatic

        Initial State: Database with multiple recordings

        Input: list\_all\_recordings() call

        Output: Sorted list of RecordingInfo objects (newest first)

        Test Case Derivation: Verifies FR-12 (recording listing)

        How test will be performed: Automated via
        test\_database.py::TestListAllRecordings

  \item{\tid{test-database-rename-delete}\\}

        Type: Functional, Dynamic, Automatic

        Initial State: Database with existing recording

        Input: rename\_recording() and delete\_recording() calls

        Output: meta.json updated with new name; directory removed on delete

        Test Case Derivation: Verifies FR-12 (recording renaming and deletion)

        How test will be performed: Automated via
        test\_database.py::TestRenameRecording, TestDeleteRecording

  \item{\tid{test-database-path-traversal-protection}\\}

        Type: Functional, Dynamic, Automatic

        Initial State: Database with validation enabled

        Input: Recording IDs containing path traversal patterns (../etc/passwd)

        Output: ValueError raised; no filesystem access outside base path

        Test Case Derivation: Verifies security requirement for path validation

        How test will be performed: Automated via test\_database.py::TestValidateId

  \item{\tid{test-database-space-usage}\\}

        Type: Functional, Dynamic, Automatic

        Initial State: Mocked disk usage

        Input: space\_usage\_bytes() and recording\_duration\_left\_s() calls

        Output: Correct used/total bytes; estimated recording duration based on free
        space

        Test Case Derivation: Verifies FR-10 (storage status reporting)

        How test will be performed: Automated via test\_database.py::TestSpaceUsage

\end{enumerate}

\subsubsection{Module: Gimbal Abstraction (\mref{mGimbal})}

The Gimbal Abstraction module handles serial communication with the gimbal
hardware.

\begin{enumerate}

  \item{\tid{test-gimbal-crc8}\\}

        Type: Functional, Dynamic, Automatic

        Initial State: GimbalSerial instance

        Input: Various byte sequences for CRC calculation

        Output: Correct CRC-8 values; consistent results; valid byte range (0-255)

        Test Case Derivation: Verifies communication protocol integrity

        How test will be performed: Automated via test\_gimbal.py::TestCrc8Smbus

  \item{\tid{test-gimbal-packet-construction}\\}

        Type: Functional, Dynamic, Automatic

        Initial State: GimbalSerial with mocked serial port

        Input: create\_request\_data() with various request IDs and payloads

        Output: Properly formatted packets with CRC at byte 0, request ID at byte 1

        Test Case Derivation: Verifies INT-2 (gimbal protocol implementation)

        How test will be performed: Automated via
        test\_gimbal.py::TestCreateRequestData

  \item{\tid{test-gimbal-angle-commands}\\}

        Type: Functional, Dynamic, Automatic

        Initial State: GimbalSerial with mocked serial port returning ACK

        Input: set\_deg(tilt, pan) and get\_deg() calls

        Output: True on success; correct angle values parsed from response

        Test Case Derivation: Verifies FR-2, FR-6, FR-7 (gimbal angle control)

        How test will be performed: Automated via test\_gimbal.py::TestSetDeg,
        TestGetDeg

  \item{\tid{test-gimbal-led-control}\\}

        Type: Functional, Dynamic, Automatic

        Initial State: GimbalSerial with mocked serial port

        Input: set\_arm\_led(True/False) and set\_status\_led(True/False)

        Output: True on success; serial write called with correct packets

        Test Case Derivation: Verifies visual feedback control for armed status

        How test will be performed: Automated via test\_gimbal.py::TestLedCommands

  \item{\tid{test-gimbal-info-query}\\}

        Type: Functional, Dynamic, Automatic

        Initial State: GimbalSerial with mocked response

        Input: gimbal\_info() call

        Output: Tilt range, pan range, and focal length range parsed correctly

        Test Case Derivation: Verifies gimbal capability detection

        How test will be performed: Automated via test\_gimbal.py::TestGimbalInfo

\end{enumerate}

\subsubsection{Module: Tracking (\mref{mTrack})}

The Tracking module implements the control law for automatic rocket tracking.

\begin{enumerate}

  \item{\tid{test-tracking-queue-behavior}\\}

        Type: Functional, Dynamic, Automatic

        Initial State: Tracking instance with mocked gimbal

        Input: Sequential on\_detection() calls with bounding box coordinates

        Output: Detections queued; oldest dropped when queue full; None accepted

        Test Case Derivation: Verifies FR-3, FR-7 (tracking queue management)

        How test will be performed: Automated via test\_tracking.py::TestOnDetection

  \item{\tid{test-tracking-worker-control-law}\\}

        Type: Functional, Dynamic, Automatic

        Initial State: Tracking instance with mocked gimbal at known position

        Input: Detection coordinates off-center from frame

        Output: gimbal.set\_deg() called with corrected angles based on proportional
        control

        Test Case Derivation: Verifies FR-7 (tracking control law)

        How test will be performed: Automated via
        test\_tracking.py::TestWorkerControlLaw

  \item{\tid{test-tracking-angle-clamping}\\}

        Type: Functional, Dynamic, Automatic

        Initial State: Tracking instance with mocked gimbal near angle limits

        Input: Detection that would push angles beyond valid range

        Output: New angles clamped to valid tilt/pan ranges

        Test Case Derivation: Verifies IR-1 (tracking command validation)

        How test will be performed: Automated via
        test\_tracking.py::TestWorkerControlLaw

  \item{\tid{test-tracking-stop}\\}

        Type: Functional, Dynamic, Automatic

        Initial State: Running Tracking instance

        Input: stop() call

        Output: Stop event set; worker thread joined successfully

        Test Case Derivation: Verifies clean shutdown of tracking thread

        How test will be performed: Automated via test\_tracking.py::TestStop

\end{enumerate}

\subsubsection{Module: CV Process Management (\mref{mCVManagement})}

The CV Process Management module handles the lifecycle and IPC with the CV
worker process.

\begin{enumerate}

  \item{\tid{test-cv-management-osd-send}\\}

        Type: Functional, Dynamic, Automatic

        Initial State: CVProcessManagement with open IPC connection

        Input: send\_osd\_data() with OSDData

        Output: Data sent over IPC; exceptions swallowed safely

        Test Case Derivation: Verifies FR-10 (OSD data transmission to CV process)

        How test will be performed: Automated via
        test\_cv\_process\_mgmt.py::TestSendOsdData

  \item{\tid{test-cv-management-recording-commands}\\}

        Type: Functional, Dynamic, Automatic

        Initial State: CVProcessManagement with open IPC connection

        Input: start\_recording() with RecordingInfo; stop\_recording()

        Output: Recording info sent; StopRecording sentinel sent

        Test Case Derivation: Verifies FR-11 (recording command relay to CV process)

        How test will be performed: Automated via
        test\_cv\_process\_mgmt.py::TestRecordingCommands

  \item{\tid{test-cv-management-pause-resume}\\}

        Type: Functional, Dynamic, Automatic

        Initial State: CVProcessManagement with active process

        Input: pause\_pipeline() and resume\_pipeline() calls

        Output: Process terminated on pause; resume event set on resume

        Test Case Derivation: Verifies pipeline control for download operations

        How test will be performed: Automated via
        test\_cv\_process\_mgmt.py::TestPauseResume

  \item{\tid{test-cv-management-recv-loop}\\}

        Type: Functional, Dynamic, Automatic

        Initial State: CVProcessManagement with mocked connection

        Input: Simulated CVData and PreviewData via recv loop

        Output: Callbacks dispatched correctly; connection closed on EOF

        Test Case Derivation: Verifies data flow from CV process to control process

        How test will be performed: Automated via
        test\_cv\_process\_mgmt.py::TestRecvLoop

\end{enumerate}

\subsubsection{Module: Live Video Process Management (\mref{mVideoManagement})}

The Live Video Process Management module supervises the live video output
process.

\begin{enumerate}

  \item{\tid{test-livestream-construction}\\}

        Type: Functional, Dynamic, Automatic

        Initial State: System ready for process management

        Input: LivestreamProcessManagement instantiation

        Output: Daemon thread started for process supervision

        Test Case Derivation: Verifies FR-8, FR-9 (live video process lifecycle)

        How test will be performed: Automated via test\_livestream\_mgmt.py

  \item{\tid{test-livestream-process-loop}\\}

        Type: Functional, Dynamic, Automatic

        Initial State: LivestreamProcessManagement instance

        Input: Simulated process loop iteration

        Output: subprocess.Popen called with livestream command; process monitored

        Test Case Derivation: Verifies live video process supervision

        How test will be performed: Automated via test\_livestream\_mgmt.py

\end{enumerate}

\subsubsection{Module: Transcode API (\mref{mAPI})}

The Transcode API provides endpoints for stabilized video preview and download.

\begin{enumerate}

  \item{\tid{test-transcode-routes-404}\\}

        Type: Functional, Dynamic, Automatic

        Initial State: Flask app with transcode routes registered; recording not found

        Input: GET /api/recordings/<id>/preview-stabilized and /download-stabilized

        Output: HTTP 404 when recording not found

        Test Case Derivation: Verifies FR-12 (proper error handling for missing
        recordings)

        How test will be performed: Automated via
        test\_transcode\_api.py::TestTranscodeRoutes

  \item{\tid{test-transcode-streaming}\\}

        Type: Functional, Dynamic, Automatic

        Initial State: Flask app with valid recording

        Input: GET preview-stabilized or download-stabilized with valid recording ID

        Output: HTTP 200 with video/webm content type; on\_download\_start/end called

        Test Case Derivation: Verifies FR-12 (stabilized video streaming)

        How test will be performed: Automated via
        test\_transcode\_api.py::TestTranscodeRoutes

  \item{\tid{test-transcode-pipe-cleanup}\\}

        Type: Functional, Dynamic, Automatic

        Initial State: Named pipe exists on filesystem

        Input: \_cleanup\_pipe() call

        Output: Pipe file removed; no error if pipe missing; OSError swallowed

        Test Case Derivation: Verifies resource cleanup for transcode operations

        How test will be performed: Automated via
        test\_transcode\_api.py::TestCleanupPipe

\end{enumerate}

\subsubsection{Module: Transcode Process (\mref{mTranscode})}

The Transcode Process module handles offline video transcoding with
stabilization.

\begin{enumerate}

  \item{\tid{test-transcode-read-log}\\}

        Type: Functional, Dynamic, Automatic

        Initial State: TranscodeProcess instance with mocked GStreamer

        Input: JSONL log file with OSD entries

        Output: List of OSDData objects parsed correctly; invalid lines skipped

        Test Case Derivation: Verifies FR-12 (log parsing for overlay generation)

        How test will be performed: Automated via
        test\_transcode\_main.py::TestTranscodeProcessReadLog

  \item{\tid{test-transcode-shader-probe}\\}

        Type: Functional, Dynamic, Automatic

        Initial State: TranscodeProcess with OSD data list

        Input: Shader probe calls with various PTS values

        Output: Correct OSD data selected; pointer advanced appropriately for preview
        vs download modes

        Test Case Derivation: Verifies overlay synchronization with video frames

        How test will be performed: Automated via
        test\_transcode\_main.py::TestTranscodeProcessShaderProbe

  \item{\tid{test-transcode-bus-call}\\}

        Type: Functional, Dynamic, Automatic

        Initial State: TranscodeProcess with mocked GStreamer bus

        Input: EOS, WARNING, and ERROR messages

        Output: Loop quit on EOS and ERROR; continued on WARNING

        Test Case Derivation: Verifies proper GStreamer pipeline event handling

        How test will be performed: Automated via
        test\_transcode\_main.py::TestTranscodeProcessBusCall

\end{enumerate}

\subsubsection{Module: CV Process (\mref{mCV})}

The CV Process module handles video acquisition, object detection, and
recording.

\begin{enumerate}

  \item{\tid{test-cv-format-time}\\}

        Type: Functional, Dynamic, Automatic

        Initial State: CV process utilities loaded

        Input: Millisecond timestamps

        Output: Human-readable formatted time strings with milliseconds

        Test Case Derivation: Verifies utility function for timestamp formatting in
        overlays

        How test will be performed: Automated via test\_cv\_main.py::TestFormatTime

  \item{\tid{test-cv-update-osd}\\}

        Type: Functional, Dynamic, Automatic

        Initial State: Mocked GStreamer osd and shader elements

        Input: OSDData with gimbal angles, GPS coordinates, IP addresses

        Output: Text property set with formatted OSD; shader uniforms updated

        Test Case Derivation: Verifies FR-8, FR-9, FR-10 (OSD overlay generation)

        How test will be performed: Automated via test\_cv\_main.py::TestUpdateOsd

\end{enumerate}

\subsubsection{Module: Common IPC (Shared Modules)}

The Common IPC module provides data structures and communication utilities.

\begin{enumerate}

  \item{\tid{test-bounding-box-center}\\}

        Type: Functional, Dynamic, Automatic

        Initial State: BoundingBox instance

        Input: Various bounding box coordinates

        Output: Correct center point (cx, cy) calculated

        Test Case Derivation: Verifies geometry calculations for tracking

        How test will be performed: Automated via test\_ipc.py::TestBoundingBoxCenter

  \item{\tid{test-bounding-box-rotate-90}\\}

        Type: Functional, Dynamic, Automatic

        Initial State: BoundingBox instance

        Input: Bounding box for 90-degree rotated frame

        Output: Correctly transformed bounding box

        Test Case Derivation: Verifies coordinate transformation for rotated camera

        How test will be performed: Automated via test\_ipc.py::TestBoundingBoxRotate90

  \item{\tid{test-dataclass-construction}\\}

        Type: Functional, Dynamic, Automatic

        Initial State: IPC module loaded

        Input: Construction of CVData, OSDData, PreviewData, RecordingInfo,
        StopRecording

        Output: Objects created with correct field values

        Test Case Derivation: Verifies data structure integrity

        How test will be performed: Automated via test\_ipc.py::TestDataclasses

\end{enumerate}

\subsubsection{Module: Common IPC Buffer}

The IPC Buffer module provides shared memory circular buffer for
high-throughput video data.

\begin{enumerate}

  \item{\tid{test-ipc-buffer-sender}\\}

        Type: Functional, Dynamic, Automatic

        Initial State: IPCBufferSender with mocked shared memory

        Input: send() calls with fixed-size payloads

        Output: Head advanced; data written at correct offset; returns False when full

        Test Case Derivation: Verifies efficient video frame buffering

        How test will be performed: Automated via
        test\_ipc\_buffer.py::TestIPCBufferSender

  \item{\tid{test-ipc-buffer-receiver}\\}

        Type: Functional, Dynamic, Automatic

        Initial State: IPCBufferReceiver with mocked shared memory containing data

        Input: receive() calls with various head/tail positions

        Output: Data retrieved correctly; tail advanced; overrun handled gracefully

        Test Case Derivation: Verifies reliable video frame retrieval

        How test will be performed: Automated via
        test\_ipc\_buffer.py::TestIPCBufferReceiver

\end{enumerate}

\subsubsection{Module: Common Watchdog}

The Watchdog module provides timeout-based process monitoring.

\begin{enumerate}

  \item{\tid{test-watchdog-refresh}\\}

        Type: Functional, Dynamic, Automatic

        Initial State: Watchdog instance with timeout and callback

        Input: refresh() calls

        Output: Timer created and started; previous timer cancelled

        Test Case Derivation: Verifies RFR-2 (watchdog functionality for process
        monitoring)

        How test will be performed: Automated via
        test\_watchdog.py::TestWatchdogRefresh

  \item{\tid{test-watchdog-clear}\\}

        Type: Functional, Dynamic, Automatic

        Initial State: Watchdog with active timer

        Input: clear() call

        Output: Timer cancelled; reference cleared; idempotent when no timer

        Test Case Derivation: Verifies safe watchdog shutdown

        How test will be performed: Automated via test\_watchdog.py::TestWatchdogClear

\end{enumerate}

\subsubsection{Module: Preview (\mref{mPreview})}

The Preview module displays live video feed in the UI.

\begin{enumerate}

  \item{\tid{test-control-page-render}\\}

        Type: Functional, Dynamic, Automatic

        Initial State: React component with mocked dependencies

        Input: Various status states (null, with/without preview, armed/disarmed,
        recording)

        Output: Preview image rendered when available; spinner shown when loading;
        REC/ARMED indicators displayed correctly

        Test Case Derivation: Verifies FR-10 (real-time preview display)

        How test will be performed: Automated via test\_pages/control.test.tsx

  \item{\tid{test-control-page-bbox-overlay}\\}

        Type: Functional, Dynamic, Automatic

        Initial State: Control page with status containing bbox

        Input: Status with bounding box coordinates and confidence

        Output: Bounding box overlay rendered with confidence value

        Test Case Derivation: Verifies FR-3 (detection visualization)

        How test will be performed: Automated via test\_pages/control.test.tsx

\end{enumerate}

\subsubsection{Module: Manual Control (\mref{mManual})}

The Manual Control module provides UI for gimbal control.

\begin{enumerate}

  \item{\tid{test-controls-card-render}\\}

        Type: Functional, Dynamic, Automatic

        Initial State: ControlsCard component with mocked API

        Input: Various system states (armed/disarmed)

        Output: Directional controls rendered; manual move API called on interaction

        Test Case Derivation: Verifies FR-2, FR-6 (manual gimbal control UI)

        How test will be performed: Automated via
        test\_components/ControlsCard.test.tsx

  \item{\tid{test-control-page-drag}\\}

        Type: Functional, Dynamic, Automatic

        Initial State: Control page in disarmed mode

        Input: Pointer drag gestures on preview area

        Output: manualMoveTo API called with calculated tilt/pan values

        Test Case Derivation: Verifies FR-2 (intuitive manual control)

        How test will be performed: Automated via test\_pages/control.test.tsx

\end{enumerate}

\subsubsection{Module: Recording Management (\mref{mRecordMng})}

The Recording Management module provides UI for listing, renaming, deleting,
and previewing recordings.

\begin{enumerate}

  \item{\tid{test-recordings-page-loading}\\}

        Type: Functional, Dynamic, Automatic

        Initial State: Recordings page mounted

        Input: API loading state

        Output: Loading spinner displayed

        Test Case Derivation: Verifies UI feedback during data fetch

        How test will be performed: Automated via test\_pages/recordings.test.tsx

  \item{\tid{test-recordings-page-empty}\\}

        Type: Functional, Dynamic, Automatic

        Initial State: Recordings page with empty list

        Input: API returns empty recordings array

        Output: ``No recordings'' message displayed

        Test Case Derivation: Verifies FR-12 (empty state handling)

        How test will be performed: Automated via test\_pages/recordings.test.tsx

  \item{\tid{test-recordings-page-list}\\}

        Type: Functional, Dynamic, Automatic

        Initial State: Recordings page with sample data

        Input: API returns list of recordings

        Output: Recording names rendered; dates and durations formatted

        Test Case Derivation: Verifies FR-12 (recording listing)

        How test will be performed: Automated via test\_pages/recordings.test.tsx

  \item{\tid{test-recordings-delete}\\}

        Type: Functional, Dynamic, Automatic

        Initial State: Recordings page with sample recording

        Input: Delete button clicked with confirmation

        Output: deleteRecording API called with correct ID

        Test Case Derivation: Verifies FR-12 (recording deletion with confirmation)

        How test will be performed: Automated via test\_pages/recordings.test.tsx

  \item{\tid{test-recordings-rename}\\}

        Type: Functional, Dynamic, Automatic

        Initial State: Recordings page with editable name

        Input: Name changed and blurred (or Escape pressed)

        Output: renameRecording API called with new name; cancelled on Escape

        Test Case Derivation: Verifies FR-12 (recording renaming)

        How test will be performed: Automated via test\_pages/recordings.test.tsx

\end{enumerate}

\subsubsection{Module: Configuration (\mref{mConfig})}

The Configuration module provides UI for system settings.

\begin{enumerate}

  \item{\tid{test-language-selector}\\}

        Type: Functional, Dynamic, Automatic

        Initial State: LanguageSelector component

        Input: Language selection change

        Output: Language atom updated; i18n activated

        Test Case Derivation: Verifies PI-2, EZ-4 (language configuration)

        How test will be performed: Automated via
        test\_components/LanguageSelector.test.tsx

  \item{\tid{test-configuration-menu}\\}

        Type: Functional, Dynamic, Automatic

        Initial State: ConfigurationMenu component

        Input: Settings changes (temperature unit, drag sensitivity, invert drag)

        Output: Setting atoms updated with new values

        Test Case Derivation: Verifies UPR-1, UPR-3 (user preference configuration)

        How test will be performed: Automated via
        test\_components/ConfigurationMenu.test.tsx

\end{enumerate}

\subsubsection{Module: API Client (Network Layer)}

The API Client module provides HTTP communication with the backend.

\begin{enumerate}

  \item{\tid{test-api-client-urls}\\}

        Type: Functional, Dynamic, Automatic

        Initial State: ApiClient instance

        Input: Base URL configuration

        Output: Correct endpoint URLs generated for status stream, recordings, previews

        Test Case Derivation: Verifies FR-10, FR-12 (API endpoint construction)

        How test will be performed: Automated via test\_network/api.test.ts

  \item{\tid{test-api-client-automatic-discovery}\\}

        Type: Functional, Dynamic, Automatic

        Initial State: Network with multiple potential server URLs

        Input: createAutomatic() with probe URLs

        Output: Client created with first responding URL; throws if all fail

        Test Case Derivation: Verifies SCR-3 (automatic server discovery)

        How test will be performed: Automated via test\_network/api.test.ts

  \item{\tid{test-api-client-methods}\\}

        Type: Functional, Dynamic, Automatic

        Initial State: ApiClient with mocked fetch

        Input: manualMove(), arm(), disarm(), startRecording(), stopRecording(),
        listRecordings(), renameRecording(), deleteRecording()

        Output: Correct HTTP methods and URLs; proper request bodies; ApiError thrown
        on failure

        Test Case Derivation: Verifies FR-2, FR-6, FR-10, FR-11, FR-12 (API operations)

        How test will be performed: Automated via test\_network/api.test.ts

\end{enumerate}

\subsubsection{Module: Rocam Provider (React Context)}

The Rocam Provider module manages API client instance and global state for
React components.

\begin{enumerate}

  \item{\tid{test-rocam-provider-initialization}\\}

        Type: Functional, Dynamic, Automatic

        Initial State: RocamProvider component with mocked createAutomatic

        Input: Provider mounted with children

        Output: API client created; status polling started; loading state managed

        Test Case Derivation: Verifies FR-10 (provider initializes API connection on
        mount)

        How test will be performed: Automated via test\_network/rocamProvider.test.tsx

  \item{\tid{test-rocam-provider-error-handling}\\}

        Type: Nonfunctional, Robustness, Automatic

        Initial State: RocamProvider with mocked API that fails to connect

        Input: API initialization failure

        Output: Error state set; error message displayed; component doesn't crash

        Test Case Derivation: Verifies RFR-1 (graceful handling of connection failures)

        How test will be performed: Automated via test\_network/rocamProvider.test.tsx

  \item{\tid{test-rocam-provider-polling}\\}

        Type: Functional, Dynamic, Automatic

        Initial State: RocamProvider with mocked API client

        Input: Status stream events and polling intervals

        Output: Status state updated; polling continues with proper cleanup on unmount

        Test Case Derivation: Verifies FR-10 (real-time status updates via SSE)

        How test will be performed: Automated via test\_network/rocamProvider.test.tsx

  \item{\tid{test-use-rocam-hook}\\}

        Type: Functional, Dynamic, Automatic

        Initial State: Component wrapped in RocamProvider

        Input: useRocam() hook called

        Output: API client and status accessible to consuming components

        Test Case Derivation: Verifies React context API exposes required values

        How test will be performed: Automated via test\_network/rocamProvider.test.tsx

\end{enumerate}

\subsubsection{Module: Utility Formatters}

The Utility Formatters module provides data formatting functions.

\begin{enumerate}

  \item{\tid{test-formatters-degrees}\\}

        Type: Functional, Dynamic, Automatic

        Initial State: Formatter functions loaded

        Input: Various numeric values for degrees, FPS, percentages

        Output: Correctly formatted strings with appropriate precision and units

        Test Case Derivation: Verifies PI-1 (intuitive value display)

        How test will be performed: Automated via test\_utils/formatters.test.ts

  \item{\tid{test-formatters-temperature}\\}

        Type: Functional, Dynamic, Automatic

        Initial State: Formatter functions loaded

        Input: Celsius temperatures; Fahrenheit unit preference

        Output: Correct conversion and formatting with degree symbol

        Test Case Derivation: Verifies temperature display in both units

        How test will be performed: Automated via test\_utils/formatters.test.ts

  \item{\tid{test-formatters-bytes}\\}

        Type: Functional, Dynamic, Automatic

        Initial State: Formatter functions loaded

        Input: Byte counts of various magnitudes

        Output: Human-readable sizes (B, KB, MB, GB)

        Test Case Derivation: Verifies FR-10 (storage display)

        How test will be performed: Automated via test\_utils/formatters.test.ts

  \item{\tid{test-formatters-duration}\\}

        Type: Functional, Dynamic, Automatic

        Initial State: Formatter functions loaded

        Input: Millisecond durations

        Output: Formatted HH:MM:SS or relative time strings

        Test Case Derivation: Verifies FR-12 (recording duration display)

        How test will be performed: Automated via test\_utils/formatters.test.ts

\end{enumerate}

\subsubsection{Module: Internationalization (i18n)}

The Internationalization module provides multi-language support.

\begin{enumerate}

  \item{\tid{test-i18n-activation}\\}

        Type: Functional, Dynamic, Automatic

        Initial State: i18n module with mocked lingui

        Input: dynamicActivate() with locale string

        Output: Messages loaded; i18n activated with correct locale

        Test Case Derivation: Verifies PI-2 (language switching)

        How test will be performed: Automated via test\_i18n.test.ts

\end{enumerate}

\subsubsection{Module: Application Routing (App)}

The App module handles page routing and navigation.

\begin{enumerate}

  \item{\tid{test-app-routing}\\}

        Type: Functional, Dynamic, Automatic

        Initial State: App component with MemoryRouter

        Input: Navigation to / and /recordings routes

        Output: Correct page components rendered; unknown paths redirect to /

        Test Case Derivation: Verifies EZ-1 (navigation between pages)

        How test will be performed: Automated via test\_App.test.tsx

\end{enumerate}

\subsection{Tests for Nonfunctional Requirements}

Unit tests for nonfunctional requirements verify performance characteristics,
reliability, and quality attributes of individual modules. These tests focus on
measurable behaviors such as timing, resource usage, error handling, and
robustness.

\subsubsection{Module: State Management (\mref{mState})}

\begin{enumerate}

  \item{\tid{test-state-status-performance}\\}

        Type: Nonfunctional, Performance, Automatic

        Initial State: StateManagement with mocked dependencies and cached metrics

        Input: status() called repeatedly

        Output: Status returned within milliseconds; cached values returned without
        recalculation

        Test Case Derivation: Verifies SLR-2, SLR-3 (status query performance for
        real-time monitoring)

        How test will be performed: Automated via
        test\_state\_management\_class.py::TestStatus

  \item{\tid{test-state-error-handling}\\}

        Type: Nonfunctional, Robustness, Automatic

        Initial State: StateManagement with mocked components that may raise exceptions

        Input: Various operations with components configured to raise exceptions

        Output: Operations fail gracefully; no unhandled exceptions propagate

        Test Case Derivation: Verifies RFR-1 (error handling for robustness)

        How test will be performed: Automated via test\_state\_management\_class.py

\end{enumerate}

\subsubsection{Module: API Gateway (\mref{mAPI})}

\begin{enumerate}

  \item{\tid{test-api-error-responses}\\}

        Type: Nonfunctional, Robustness, Automatic

        Initial State: Flask test client

        Input: Invalid JSON, missing required fields, malformed requests

        Output: HTTP 400 with descriptive error messages; no server crashes

        Test Case Derivation: Verifies RFR-1 (API error handling and feedback)

        How test will be performed: Automated via test\_api.py::TestRecordingsRename

\end{enumerate}

\subsubsection{Module: SSE (Server-Sent Events)}

\begin{enumerate}

  \item{\tid{test-sse-broadcast-performance}\\}

        Type: Nonfunctional, Performance, Automatic

        Initial State: MessageAnnouncer with multiple listeners registered

        Input: Rapid announce() calls with status updates

        Output: All listeners receive messages; full queues drop messages silently (no
        blocking)

        Test Case Derivation: Verifies SLR-3, RFR-4 (real-time status streaming without
        backpressure)

        How test will be performed: Automated via test\_sse.py::TestMessageAnnouncer

  \item{\tid{test-sse-thread-safety}\\}

        Type: Nonfunctional, Robustness, Automatic

        Initial State: MessageAnnouncer with concurrent operations

        Input: Simultaneous listen(), remove(), and announce() from multiple threads

        Output: No race conditions; all operations complete safely

        Test Case Derivation: Verifies thread safety for concurrent UI connections

        How test will be performed: Automated via test\_sse.py::TestMessageAnnouncer

\end{enumerate}

\subsubsection{Module: IPC Buffer}

\begin{enumerate}

  \item{\tid{test-ipc-buffer-overrun-handling}\\}

        Type: Nonfunctional, Robustness, Automatic

        Initial State: IPCBufferReceiver with head significantly ahead of tail (overrun
        condition)

        Input: receive() call during overrun

        Output: Tail advanced to catch up; data retrieved without corruption

        Test Case Derivation: Verifies RFR-4 (graceful handling of producer/consumer
        rate mismatch)

        How test will be performed: Automated via
        test\_ipc\_buffer.py::TestIPCBufferReceiver

  \item{\tid{test-ipc-buffer-memory-safety}\\}

        Type: Nonfunctional, Robustness, Automatic

        Initial State: IPCBuffer with various head/tail positions including edge cases

        Input: Operations at buffer boundaries

        Output: No out-of-bounds access; proper ring buffer wrapping

        Test Case Derivation: Verifies memory safety for shared video data

        How test will be performed: Automated via test\_ipc\_buffer.py

\end{enumerate}

\subsubsection{Module: CV Process Management (\mref{mCVManagement})}

\begin{enumerate}

  \item{\tid{test-cv-management-exception-safety}\\}

        Type: Nonfunctional, Robustness, Automatic

        Initial State: CVProcessManagement with connection in various states

        Input: Operations during broken pipe, EOF, and other error conditions

        Output: Exceptions caught and swallowed; component continues operating

        Test Case Derivation: Verifies RFR-2 (fault tolerance for process
        communication)

        How test will be performed: Automated via
        test\_cv\_process\_mgmt.py::TestExceptionHandling

\end{enumerate}

\subsubsection{Module: Recording Database (\mref{mRecord})}

\begin{enumerate}

  \item{\tid{test-database-error-handling}\\}

        Type: Nonfunctional, Robustness, Automatic

        Initial State: Database with simulated filesystem errors

        Input: Operations with mocked os/path functions raising OSError

        Output: Graceful degradation; sensible defaults returned; exceptions handled

        Test Case Derivation: Verifies RFR-1, CR-1 (error handling and data integrity)

        How test will be performed: Automated via test\_database.py::TestGetSizeBytes

\end{enumerate}

\subsubsection{Module: Watchdog}

\begin{enumerate}

  \item{\tid{test-watchdog-timeout-accuracy}\\}

        Type: Nonfunctional, Performance, Automatic

        Initial State: Watchdog with specific timeout value

        Input: refresh() calls at various intervals

        Output: Timer created with correct timeout; callback triggered after timeout

        Test Case Derivation: Verifies PAR-1 (timing accuracy for process monitoring)

        How test will be performed: Automated via test\_watchdog.py

\end{enumerate}

\subsubsection{Module: System Status (\mref{mStatus})}

\begin{enumerate}

  \item{\tid{test-system-status-caching}\\}

        Type: Nonfunctional, Performance, Automatic

        Initial State: SystemStatusMonitor with mocked jtop

        Input: Multiple getter calls between updates

        Output: Cached values returned immediately; no redundant hardware queries

        Test Case Derivation: Verifies SLR-2, SLR-3 (efficient status monitoring)

        How test will be performed: Automated via test\_system\_status.py

  \item{\tid{test-system-status-callback-updates}\\}

        Type: Nonfunctional, Correctness, Automatic

        Initial State: SystemStatusMonitor with initial readings

        Input: jtop callback with updated sensor values

        Output: All cached metrics updated to new values

        Test Case Derivation: Verifies EZ-3 (accurate system telemetry)

        How test will be performed: Automated via
        test\_system\_status.py::TestSystemStatusMonitor

\end{enumerate}

\subsubsection{Module: Common Utils}

The Utils module provides network and scheduling utilities.

\begin{enumerate}

  \item{\tid{test-ip4-addresses}\\}

        Type: Functional, Dynamic, Automatic

        Initial State: Utils module with mocked netifaces

        Input: Network interface configurations with various IPv4 addresses

        Output: List of IP addresses extracted correctly; empty list when no interfaces

        Test Case Derivation: Verifies FR-10 (device IP address reporting for UI)

        How test will be performed: Automated via test\_utils.py::TestIp4Addresses

  \item{\tid{test-scheduler-validation}\\}

        Type: Nonfunctional, Robustness, Automatic

        Initial State: Scheduler functions with mocked OS calls

        Input: Various priority values (valid and invalid) for FIFO, OTHER, and BATCH
        schedulers

        Output: ValueError raised for out-of-range priorities; successful calls for
        valid values

        Test Case Derivation: Verifies PAR-1 (process scheduling priority validation)

        How test will be performed: Automated via
        test\_utils.py::TestSchedulerValidation

\end{enumerate}

\subsubsection{Module: Recording-Only State Management}

The Recording-Only State Management module provides a limited state management
implementation for recording management mode.

\begin{enumerate}

  \item{\tid{test-state-management-recording-only-noops}\\}

        Type: Functional, Dynamic, Automatic

        Initial State: StateManagementRecordingManagementOnly instance

        Input: Calls to arm(), disarm(), manual\_move(), manual\_move\_to(),
        start\_recording(), stop\_recording()

        Output: All operations are no-ops (do not raise exceptions)

        Test Case Derivation: Verifies that recording-only mode safely ignores control
        operations

        How test will be performed: Automated via
        test\_state\_management\_recording\_only.py::TestStateManagementRecordingOnly

  \item{\tid{test-state-management-recording-only-status}\\}

        Type: Functional, Dynamic, Automatic

        Initial State: StateManagementRecordingManagementOnly with mocked database and
        system status

        Input: status() call

        Output: StatusResponse with is\_recording=False, armed=False, system metrics
        populated

        Test Case Derivation: Verifies FR-10 (status reporting in recording-only mode)

        How test will be performed: Automated via
        test\_state\_management\_recording\_only.py::TestStateManagementRecordingOnly

  \item{\tid{test-control-process-main-entry}\\}

        Type: Functional, Dynamic, Automatic

        Initial State: Mocked StateManagement and threading

        Input: run\_control\_process() and run\_recording\_management() calls

        Output: API thread started and joined; StateManagement instantiated

        Test Case Derivation: Verifies control process entry points initialize
        correctly

        How test will be performed: Automated via
        test\_state\_management\_recording\_only.py::TestControlProcessMain

\end{enumerate}

\subsubsection{Module: Main Entry Point}

The Main Entry Point module handles command-line dispatch to different process
modes.

\begin{enumerate}

  \item{\tid{test-main-dispatch-control}\\}

        Type: Functional, Dynamic, Automatic

        Initial State: Mocked sys.argv with no arguments

        Input: Default execution (no command line args)

        Output: run\_control\_process() called

        Test Case Derivation: Verifies default mode launches control process

        How test will be performed: Automated via
        test\_main\_entry.py::TestMainDispatch

  \item{\tid{test-main-dispatch-cv}\\}

        Type: Functional, Dynamic, Automatic

        Initial State: Mocked sys.argv with ["cv"] argument

        Input: CV process launch command

        Output: run\_cv\_process() called

        Test Case Derivation: Verifies FR-3 (CV process launch)

        How test will be performed: Automated via
        test\_main\_entry.py::TestMainDispatch

  \item{\tid{test-main-dispatch-livestream}\\}

        Type: Functional, Dynamic, Automatic

        Initial State: Mocked sys.argv with ["livestream"] argument

        Input: Livestream process launch command

        Output: run\_livestream\_process() called

        Test Case Derivation: Verifies FR-8, FR-9 (live video process launch)

        How test will be performed: Automated via
        test\_main\_entry.py::TestMainDispatch

  \item{\tid{test-main-dispatch-transcode}\\}

        Type: Functional, Dynamic, Automatic

        Initial State: Mocked sys.argv with ["download-stabilized"] or
          ["preview-stabilized"] arguments

        Input: Transcode process launch command with video and log paths

        Output: start\_transcode\_process() called with correct mode and paths

        Test Case Derivation: Verifies FR-12 (transcode process launch for recording
        stabilization)

        How test will be performed: Automated via
        test\_main\_entry.py::TestMainDispatch

  \item{\tid{test-main-dispatch-unknown}\\}

        Type: Nonfunctional, Robustness, Automatic

        Initial State: Mocked sys.argv with unknown command

        Input: Unknown command line argument

        Output: System exits with code 1

        Test Case Derivation: Verifies proper error handling for invalid commands

        How test will be performed: Automated via
        test\_main\_entry.py::TestMainDispatch

\end{enumerate}

\subsubsection{Module: Utility Formatters}

\begin{enumerate}

  \item{\tid{test-formatters-edge-cases}\\}

        Type: Nonfunctional, Robustness, Automatic

        Initial State: Formatter functions loaded

        Input: Edge values (null, undefined, negative, zero, very large numbers)

        Output: Sensible defaults; no crashes; formatted output for all inputs

        Test Case Derivation: Verifies RFR-1 (robustness against unexpected data)

        How test will be performed: Automated via test\_utils/formatters.test.ts

  \item{\tid{test-formatters-internationalization}\\}

        Type: Nonfunctional, Internationalization, Automatic

        Initial State: Formatters with various locale considerations

        Input: Values requiring locale-specific formatting

        Output: Consistent formatting regardless of system locale

        Test Case Derivation: Verifies PI-2 (consistent international display)

        How test will be performed: Automated via test\_utils/formatters.test.ts

\end{enumerate}

\subsubsection{Module: UI Components}

\begin{enumerate}

  \item{\tid{test-ui-loading-states}\\}

        Type: Nonfunctional, Usability, Automatic

        Initial State: Components with loading/null states

        Input: Null or undefined data

        Output: Loading spinners or placeholder content; no crashes

        Test Case Derivation: Verifies EZ-2 (immediate feedback for pending operations)

        How test will be performed: Automated via component test files (e.g.,
        test\_pages/control.test.tsx)

  \item{\tid{test-ui-error-handling}\\}

        Type: Nonfunctional, Robustness, Automatic

        Initial State: Components with mocked API that returns errors

        Input: API errors, network failures, malformed responses

        Output: Graceful error display; component remains interactive

        Test Case Derivation: Verifies RFR-1 (UI error resilience)

        How test will be performed: Automated via component test files

\end{enumerate}

\subsection{Traceability Between Test Cases and Modules}

The following tables provide traceability between unit test cases and the
modules defined in the Module Guide (MG). Each test case verifies specific
functionality of its corresponding module, ensuring that all implemented leaf
modules have adequate test coverage.

\begin{longtable}{|p{4.5cm}|p{10cm}|}
  \caption{Traceability Between Unit Test Cases and Backend Modules} \label{TblUTM-Backend}                                                                                                                                                                                                                                                      \\
  \hline
  \textbf{Module}                                    & \textbf{Unit Test Cases}                                                                                                                                                                                                                                                                  \\
  \hline
  \endfirsthead
  \multicolumn{2}{c}{\tablename\ \thetable{} -- Continued from previous page}                                                                                                                                                                                                                                                                    \\
  \hline
  \textbf{Module}                                    & \textbf{Unit Test Cases}                                                                                                                                                                                                                                                                  \\
  \hline
  \endhead
  \hline
  \multicolumn{2}{r}{Continued on next page}                                                                                                                                                                                                                                                                                                     \\
  \endfoot
  \hline
  \endlastfoot
  \mref{mState} (State Management)                   & \tid{test-state-management-arm-disarm}; \tid{test-state-management-manual-move}; \tid{test-state-management-manual-move-clamped}; \tid{test-state-management-recording}; \tid{test-state-management-download-pause}; \tid{test-state-status-performance}; \tid{test-state-error-handling} \\
  \hline
  \mref{mAPI} (API Gateway)                          & \tid{test-api-endpoints-exist}; \tid{test-api-recordings-crud}; \tid{test-api-validation}; \tid{test-api-error-responses}                                                                                                                                                                 \\
  \hline
  \mref{mRecord} (Recording Database)                & \tid{test-database-allocate}; \tid{test-database-list-recordings}; \tid{test-database-rename-delete}; \tid{test-database-path-traversal-protection}; \tid{test-database-space-usage}; \tid{test-database-error-handling}                                                                  \\
  \hline
  \mref{mGimbal} (Gimbal Abstraction)                & \tid{test-gimbal-crc8}; \tid{test-gimbal-packet-construction}; \tid{test-gimbal-angle-commands}; \tid{test-gimbal-led-control}; \tid{test-gimbal-info-query}                                                                                                                              \\
  \hline
  \mref{mTrack} (Tracking)                           & \tid{test-tracking-queue-behavior}; \tid{test-tracking-worker-control-law}; \tid{test-tracking-angle-clamping}; \tid{test-tracking-stop}                                                                                                                                                  \\
  \hline
  \mref{mCVManagement} (CV Process Management)       & \tid{test-cv-management-osd-send}; \tid{test-cv-management-recording-commands}; \tid{test-cv-management-pause-resume}; \tid{test-cv-management-recv-loop}; \tid{test-cv-management-exception-safety}                                                                                      \\
  \hline
  \mref{mVideoManagement} (Live Video Management)    & \tid{test-livestream-construction}; \tid{test-livestream-process-loop}                                                                                                                                                                                                                    \\
  \hline
  \mref{mTranscodeManagement} (Transcode Management) & \tid{test-transcode-read-log}; \tid{test-transcode-shader-probe}; \tid{test-transcode-bus-call}                                                                                                                                                                                           \\
  \hline
  \mref{mTranscode} (Transcode Process)              & \tid{test-transcode-read-log}; \tid{test-transcode-shader-probe}; \tid{test-transcode-bus-call}                                                                                                                                                                                           \\
  \hline
  \mref{mCV} (CV Process)                            & \tid{test-cv-format-time}; \tid{test-cv-update-osd}                                                                                                                                                                                                                                       \\
  \hline
  \mref{mVideo} (Live Video Process)                 & \tid{test-livestream-construction}; \tid{test-livestream-process-loop}                                                                                                                                                                                                                    \\
  \hline
  Common IPC                                         & \tid{test-bounding-box-center}; \tid{test-bounding-box-rotate-90}; \tid{test-dataclass-construction}                                                                                                                                                                                      \\
  \hline
  Common IPC Buffer                                  & \tid{test-ipc-buffer-sender}; \tid{test-ipc-buffer-receiver}; \tid{test-ipc-buffer-overrun-handling}; \tid{test-ipc-buffer-memory-safety}                                                                                                                                                 \\
  \hline
  Common Watchdog                                    & \tid{test-watchdog-refresh}; \tid{test-watchdog-clear}; \tid{test-watchdog-timeout-accuracy}                                                                                                                                                                                              \\
  \hline
  Common System Status                               & \tid{test-system-status-caching}; \tid{test-system-status-callback-updates}                                                                                                                                                                                                               \\
  \hline
  Common Utils                                       & \tid{test-ip4-addresses}; \tid{test-scheduler-validation}                                                                                                                                                                                                                                 \\
  \hline
  Recording-Only State Management                    & \tid{test-state-management-recording-only-noops}; \tid{test-state-management-recording-only-status}; \tid{test-control-process-main-entry}                                                                                                                                                \\
  \hline
  Main Entry Point                                   & \tid{test-main-dispatch-control}; \tid{test-main-dispatch-cv}; \tid{test-main-dispatch-livestream}; \tid{test-main-dispatch-transcode}; \tid{test-main-dispatch-unknown}                                                                                                                  \\
  \hline
  Transcode API (\mref{mAPI})                        & \tid{test-transcode-routes-404}; \tid{test-transcode-streaming}; \tid{test-transcode-pipe-cleanup}                                                                                                                                                                                        \\
  \hline
  SSE (\mref{mAPI})                                  & \tid{test-sse-broadcast-performance}; \tid{test-sse-thread-safety}                                                                                                                                                                                                                        \\
\end{longtable}

\begin{table}[H]
  \centering
  \caption{Traceability Between Unit Test Cases and Frontend Modules}
  \label{TblUTM3}
  \begin{tabular}{|p{4.5cm}|p{10cm}|}
    \hline
    \textbf{Module}                          & \textbf{Unit Test Cases}                                                                                                                                                                                    \\
    \hline
    \mref{mPreview} (Preview)                & \tid{test-control-page-render}; \tid{test-control-page-bbox-overlay}                                                                                                                                        \\
    \hline
    \mref{mManual} (Manual Control)          & \tid{test-controls-card-render}; \tid{test-control-page-drag}                                                                                                                                               \\
    \hline
    \mref{mRecordMng} (Recording Management) & \tid{test-recordings-page-loading}; \tid{test-recordings-page-empty}; \tid{test-recordings-page-list}; \tid{test-recordings-delete}; \tid{test-recordings-rename}                                           \\
    \hline
    \mref{mConfig} (Configuration)           & \tid{test-language-selector}; \tid{test-configuration-menu}                                                                                                                                                 \\
    \hline
    API Client                               & \tid{test-api-client-urls}; \tid{test-api-client-automatic-discovery}; \tid{test-api-client-methods}                                                                                                        \\
    \hline
    Rocam Provider                           & \tid{test-rocam-provider-initialization}; \tid{test-rocam-provider-error-handling}; \tid{test-rocam-provider-polling}; \tid{test-use-rocam-hook}                                                            \\
    \hline
    Utility Formatters                       & \tid{test-formatters-degrees}; \tid{test-formatters-temperature}; \tid{test-formatters-bytes}; \tid{test-formatters-duration}; \tid{test-formatters-edge-cases}; \tid{test-formatters-internationalization} \\
    \hline
    Internationalization                     & \tid{test-i18n-activation}                                                                                                                                                                                  \\
    \hline
    App (Routing)                            & \tid{test-app-routing}                                                                                                                                                                                      \\
    \hline
    UI Components (Individual)               & \tid{test-system-status-card}; \tid{test-navbar}; \tid{test-controls-card}; \tid{test-language-selector}; \tid{test-configuration-menu}; \tid{test-ui-loading-states}; \tid{test-ui-error-handling}         \\
    \hline
  \end{tabular}
\end{table}

\textbf{Note:} The following modules are listed in MG but are not verified by unit tests (as documented in the Unit Testing Scope section):
\begin{itemize}
  \item \mref{mStream} (Video Stream Abstraction) --- Third-party library (Nvidia Deepstream SDK)
  \item \mref{mSerial} (Serial Abstraction) --- Third-party library (pyserial)
  \item \mref{mStatus} (System Status Module) --- Hardware-dependent library (jetson-stats)
\end{itemize}

Container modules (\mref{mJetson}, \mref{mControl}, \mref{mUI}) are also not
directly tested as they are organizational constructs; their leaf modules are
tested individually.

\bibliographystyle{plainnat}

\bibliography{../../refs/References}

\newpage

\section{Appendix}

\subsection{User Survey Questions}

Please answer each question based on your experience during the field test. For
questions rated on a scale, circle the number that best reflects your opinion:
1 = Strongly Disagree, 2 = Disagree, 3 = Neutral, 4 = Agree, 5 = Strongly
Agree.

\subsubsection*{User Interface}

\begin{enumerate}
  \item The interface was readable and easy to see outdoors. \hfill (1 2 3 4 5)
  \item The interface looked professional and well-designed. \hfill (1 2 3 4 5)
  \item The important information was immediately visible on the screen. \hfill (1 2 3
        4 5)
  \item No unnecessary or confusing information was shown. \hfill (1 2 3 4 5)
  \item The system was easy to learn and operate after a short time. \hfill (1 2 3 4 5)
  \item The terms and labels used in the interface were clear and understandable.
        \hfill (1 2 3 4 5)
\end{enumerate}

\subsubsection*{System Performance}

\begin{enumerate}
  \item The tracking footage was stable and smooth. \hfill (1 2 3 4 5)
  \item The tracking system maintained lock on the rocket reliably. \hfill (1 2 3 4 5)
  \item The system operated as expected without major errors. \hfill (1 2 3 4 5)
  \item Once armed, the system operated fully autonomously. \hfill (1 2 3 4 5)
  \item I was able to fully operate and monitor the system remotely. \hfill (1 2 3 4 5)
\end{enumerate}

\subsubsection*{User Experience and Feedback}

\begin{enumerate}
  \item Were any colors, symbols, or content elements disrespectful or inappropriate to
        you?\\ \textit{(Please describe if applicable)}\\[1em]
        \makebox[0.95\textwidth]{\hrulefill}\\[0.5em]
        \makebox[0.95\textwidth]{\hrulefill}
  \item What improvements would you suggest for future versions of RoCam?\\[1em]
        \makebox[0.95\textwidth]{\hrulefill}\\[0.5em]
        \makebox[0.95\textwidth]{\hrulefill}
\end{enumerate}

\newpage{}
\section*{Appendix --- Reflection}

The information in this section will be used to evaluate the team members on
the graduate attribute of Lifelong Learning.

\input{../Reflection.tex}

\begin{enumerate}
  \item What went well while writing this deliverable?
  \item What pain points did you experience during this deliverable, and how did you
        resolve them?
  \item What knowledge and skills will the team collectively need to acquire to
        successfully complete the verification and validation of your project? Examples
        of possible knowledge and skills include dynamic testing knowledge, static
        testing knowledge, specific tool usage, Valgrind etc. You should look to
        identify at least one item for each team member.
  \item For each of the knowledge areas and skills identified in the previous question,
        what are at least two approaches to acquiring the knowledge or mastering the
        skill? Of the identified approaches, which will each team member pursue, and
        why did they make this choice?
\end{enumerate}

\subsection*{Jianqing Liu}
\begin{enumerate}
  \item \textbf{What went well:}
        I was responsible for the Part~4 of this deliverable, focusing on all the system tests. I think it
        was really easy to come up with all the test cases because We already completed the SRS.

  \item \textbf{Pain points and resolution:}
        I think some of the test cases are very verbose. For some of the tests, how test will be performed
        is just the same as the input and output condition of the test case.
\end{enumerate}

\subsection*{Mike Chen}
\begin{enumerate}
  \item \textbf{What went well:}
        I contributed primarily to Part~3 of this deliverable, focusing on the verification and validation plan
        and the automation tools section. The structure and clarity of this part improved substantially after
        integrating team and stakeholder feedback. I ensured that each verification activity was directly tied
        to the SRS and our test framework.

  \item \textbf{Pain points and resolution:}
        I initially faced several issues with GitHub CI configuration—some commits overwrote teammates’
        work and caused LaTeX compilation failures. I resolved these by restoring the affected commits,
        redesigning the YAML workflows, and setting up branch protection rules. Revisions based on
        feedback also required extensive reformatting to achieve consistency across sections.

  \item \textbf{Knowledge and skills to acquire:}
        I plan to improve my skills in CI/CD automation and verification planning, particularly integrating
        automated builds, test execution, and documentation generation through GitHub Actions.

  \item \textbf{Approaches to acquire skills:}
        I will (1) study advanced GitHub Actions pipelines from open-source projects and (2) create a
        smaller test repository to simulate multi-stage verification workflows. I selected the second
        approach because it provides hands-on experience with merge control, build automation, and
        verification feedback integration.
\end{enumerate}

\subsection*{Frank}
\begin{enumerate}
  \item \textbf{What went well:}
        I contributed to both Part~2 and Part~3 of the deliverable, ensuring that the design and verification
        sections were fully aligned with the SRS. My work focused on refining the test structure, verifying
        requirement traceability, and improving the flow between verification items and their
        corresponding system requirements.

  \item \textbf{Pain points and resolution:}
        The main challenge was maintaining consistent terminology and depth of explanation between
        sections written by different team members. To resolve this, I reviewed the SRS and cross-checked
        every requirement reference to ensure correctness and coherence across parts.

  \item \textbf{Knowledge and skills to acquire:}
        I plan to strengthen my knowledge of requirements traceability and automated testing frameworks
        to improve coverage and maintain consistency with system goals.

  \item \textbf{Approaches to acquire skills:}
        I will (1) review documentation on structured traceability mapping and (2) develop example UI
        and system tests that link directly to specific SRS requirements. I chose the second approach
        because applying traceability in practice will reinforce both documentation and test alignment.
\end{enumerate}

\subsection*{Xiaotian Lou}
\begin{enumerate}
  \item \textbf{What went well:}
        I focused mainly on Part~2 of the deliverable, organizing the SRS and design verification sections.
        The verification checklists and phase structure improved document readability and made the review
        process more systematic.

  \item \textbf{Pain points and resolution:}
        The primary difficulty was harmonizing writing styles and section formatting across contributors.
        I standardized tables, checklists, and layout conventions to produce a consistent and professional
        presentation.

  \item \textbf{Knowledge and skills to acquire:}
        I aim to gain more experience in formal verification documentation and structured validation
        reporting to improve the traceability and completeness of future deliverables.

  \item \textbf{Approaches to acquire skills:}
        I will (1) study professional verification and validation templates from engineering projects and
        (2) apply checklist-based verification methods in upcoming development stages. I selected the
        first approach to strengthen my understanding of documentation standards and best practices.
\end{enumerate}

\end{document}

